% !TeX program = xelatex
% Study notes with restrained red accents. Build: python build.py notes-starter
\AddToHook{package/projlib-language/before}{\RequirePackage{fvextra}}
\documentclass[11pt,a4paper,plain,CN,colored-proof=false,
  load-custom-latin-font-file=pku-font-latin.tex,
  load-custom-math-font-file=pku-font-math.tex,
  load-custom-cjk-font-file=pku-font-cjk.tex]{einfart}
\input{settings}
\renewcommand{\PKUDocumentType}{notes}
\renewcommand{\PKUShortTitle}{应用线性回归 R 实操笔记}
\renewcommand{\PKUSubtitle}{从方法到代码：四讲课件、课堂代码与四次作业的 R 落地}
\usepackage{pku-notes}
% R 输出框：灰字、无行号，与代码框区分
\newtcblisting{ROut}[1][输出]{pkucode,listing only,colback=white,
  colbacktitle=white,fonttitle=\sffamily\small\color{black!60},
  listing options={basicstyle=\ttfamily\footnotesize\color{black!78},
    columns=fullflexible,keepspaces=true,breaklines=true,numbers=none,
    upquote=true,showstringspaces=false,extendedchars=false},title={#1}}
\newcommand{\fn}[1]{\texttt{#1}}
% 行内代码允许在任意位置断行，避免长代码撑出版心
\fvinlineset{breaklines=true,breakanywhere=true}
% R 代码行常带对齐注释，代码框字号小一档，减少折行
\lstdefinestyle{pkuinline}{style=pkuinlinebase,basicstyle=\ttfamily\footnotesize}
\title{应用线性回归 R 实操笔记}
\author{学习网站：\url{https://fxt-gw-pb.github.io/public-health-statistics-notes/}\\[2pt]2026 年 11 月 19 日}
\date{}

\begin{document}
\MakeNotesTitle
\begin{notesabstract}
课件笔记已经把模型、估计、检验与诊断的理论讲清楚了，这份笔记只做一件事：把每一种方法落到 R 代码上，并说明输出里该读哪几个数、怎样写进结果、哪里最容易出错。内容按四讲课件的顺序展开，第 1 章补齐 R 的数据处理基础，第 2–8 章依次对应多重回归、变量筛选、回归诊断、共线性与岭回归、一般线性模型、协方差分析和重复测量方差分析，第 9 章逐份解读四次作业的代码。所有代码都来自同目录下可直接运行的脚本，文中给出的数值都是在 R 4.5.2 上实际运行得到的；与课件或课堂代码不一致的地方，在相应位置作了说明。
\end{notesabstract}

\tableofcontents
\setcounter{section}{-1}
\section{学习包的结构与用法}

\subsection{文件夹}
学习包的文件夹名为 \Verb|应用线性回归_R实操/|。在本机，它放在 \Verb|应用线性回归/| 下，与原有的 \Verb|回归1|–\Verb|回归4|、\Verb|hw1|–\Verb|hw4| 并列，原文件夹没有改动；从学习网站下载的压缩包解压后也是这个文件夹。结构如下：

\begin{ROut}[目录结构]
应用线性回归_R实操/
|-- 应用线性回归_R实操.Rproj   ← 双击它打开 RStudio
|-- 00_安装与检查.R            ← 第一次打开时运行一次
|-- README.md
|-- data/                      ← 全部数据（24 个 CSV）
|   |-- 数据说明.md            ← 每个文件的来源、变量、在哪里用到
|   `-- make_inline_data.R     ← 把课件里写在代码中的数据导出为 CSV
|-- scripts/                   ← 按讲次整理的可运行脚本（本笔记第 1–8 章、附录 B）
|   |-- 01_R工作流与数据处理.R
|   |-- 02_多重回归_拟合与解读.R
|   |-- 03_假设检验与变量筛选.R
|   |-- 04_回归诊断.R
|   |-- 05_共线性与岭回归.R
|   |-- 06_一般线性模型_分类变量与方差分析.R
|   |-- 07_协方差分析.R
|   |-- 08_重复测量方差分析.R
|   `-- 09_课件笔记例题复现.R
|-- homework/                  ← 四次作业的代码（本笔记第 9 章）
|-- 课堂原始代码/               ← 老师发的原件，未改动；附对照说明.md
|-- output/                    ← 运行脚本时生成的图和表
`-- 笔记/                      ← 本笔记的 PDF 与 LaTeX 源文件
\end{ROut}

\subsection{打开与运行}
\begin{enumerate}
\item 双击 \Verb|应用线性回归_R实操.Rproj|。RStudio 会把工作目录设为学习包根目录，所有脚本里的 \Verb|read.csv("data/...")| 都能直接找到文件。不要用“直接双击 .R 文件”的方式打开，那样工作目录不对。
\item 第一次使用先运行 \Verb|00_安装与检查.R|，它会装好缺少的包（car、leaps、olsrr、ppcor、lmtest、sandwich、emmeans、afex、ez、tidyr、ggplot2），并检查数据是否齐全。
\item 打开 \Verb|scripts/| 下的脚本，按 Ctrl/Cmd + Shift + O 显示右侧大纲。脚本里形如 \Verb|# 2.7 预测：均数的 CI 与个体的 PI ----| 的行就是大纲条目。脚本编号与本笔记的章号相同（脚本 02 对应第 2 章），正文提到“脚本 §2.7”时，指的就是大纲里的这一条。
\item 光标放在某一行按 Ctrl/Cmd + Enter 运行这一行（或选中的几行）；Ctrl/Cmd + Shift + Enter 从头运行整份脚本。每份脚本都已从头到尾跑通。
\end{enumerate}

\subsection{怎样配合本笔记使用}
每章开头的小字注明对应的课件页码、脚本和数据。正文只摘出关键的几行代码和输出，完整代码在脚本里。建议的顺序是：先读本章正文弄清“要做什么、读哪几个数”，再在 RStudio 里逐段运行脚本，对照正文核对数值，最后打开对应的作业脚本看同样的方法怎么用在作业上。

\begin{table}[htbp]
\centering\small
\caption{课件、脚本、数据与作业的对应关系}
\label{tab:map}
\begin{tabular}{@{}p{3.3cm}p{4.2cm}p{4.3cm}p{2.0cm}@{}}
\toprule
课件 & 脚本 & 主要数据 & 对应作业\\
\midrule
—— & 01 R 工作流与数据处理 & 各数据 & ——\\
第一讲 模型与解释 & 02 多重回归：拟合与解读 & datasetA & 作业一\\
第一讲后半 变量筛选 & 03 假设检验与变量筛选 & datasetC、diabetes27 & 作业一\\
第二讲 回归诊断 & 04 回归诊断 & datasetB & 作业二第一题\\
第二讲 共线性与岭回归 & 05 共线性与岭回归 & datasetB & 作业二第二题\\
第三讲 一般线性模型 & 06 一般线性模型 & birthwt、glucose\_block、sbp\_factorial2x2 & 作业三\\
第三、四讲 协方差分析 & 07 协方差分析 & datasetC\_glm、bp、knee & 作业四第 1 题\\
第四讲 重复测量 & 08 重复测量方差分析 & dm、df\_long & 作业四第 2 题\\
课件笔记各章例题 & 09 课件笔记例题复现 & 12 个小数据集 & ——\\
\bottomrule
\end{tabular}
\TableNote{数据文件均在 \texttt{data/} 下；文件来源与变量见 \texttt{data/数据说明.md}。}
\end{table}

\subsection{网页版与下载}
这份笔记也收进了学习网站《医学统计学 · 学习笔记》，是其中的第四门课程“应用线性回归 · R 实操”：

\begin{center}
\url{https://fxt-gw-pb.github.io/public-health-statistics-notes/#/course/rcode}
\end{center}

网页版可以在线阅读、全文搜索、一键复制代码块；课程页提供两个下载入口：本笔记的 PDF，以及完整的 R 学习包（\Verb|.Rproj|、全部数据、九份讲次脚本、四份作业脚本、课堂原始代码和本笔记的 LaTeX 源文件）。下载学习包解压后，按上一节的步骤打开即可。

\begin{notebox}[关于课堂原始代码]
老师发的脚本保存在 \Verb|课堂原始代码/| 里，没有改动。它们写的是 \Verb|read.csv("d://data/...")| 这样的绝对路径，又大量使用 \Verb|attach()|，有几处直接运行会报错（例如第四讲读入的对象叫 \Verb|da| 后面却用 \Verb|bp|）。整理版把这些问题都改掉了，逐条对照见 \Verb|课堂原始代码/对照说明.md| 和附录 D。
\end{notebox}

\section{R 工作流：从读数据到取结果}
\srcnote{脚本：\texttt{scripts/01\_R工作流与数据处理.R}\quad 数据：datasetA、bp、dm、sbp\_factorial3x3}

课上的代码默认大家已经会 R 的基本操作，但作业里卡住的地方，多半不是模型本身，而是数据没读对、因子水平顺序不对、长宽格式不对，或者不知道结果存在模型对象的哪里。这一章把这些操作集中过一遍。

\subsection{路径与读入}
在 \Verb|.Rproj| 打开的项目里，工作目录就是项目根目录，路径一律写相对路径：

\begin{RCode}[R · 读入数据]
getwd()                                   # 应当是 .../应用线性回归_R实操
d  <- read.csv("data/datasetA.csv")       # 默认 header = TRUE，逗号分隔
bp <- read.csv("data/bp.csv", fileEncoding = "UTF-8")   # 含中文的文件加编码
f3 <- read.csv("data/sbp_factorial3x3.csv", check.names = FALSE,
               fileEncoding = "UTF-8")    # 中文列名原样保留
names(f3) <- c("id", "A", "B", "y")       # 再改成短的英文名
read.csv("data/datasetC.csv")[, -1]       # [, -1]：去掉第 1 列 id
\end{RCode}

几个参数的作用：\Verb|fileEncoding = "UTF-8"| 防止中文乱码，在 Windows 上尤其需要；\Verb|check.names = FALSE| 让列名原样保留（默认会改写不合规的列名，例如以数字开头的列名前面加 \Verb|X|，带 BOM 的首列变成 \Verb|X.U.FEFF.编号|）；课堂代码里常见的 \Verb|[, -1]| 是把第一列（id 或行号）去掉，只留分析变量，这样后面可以写 \Verb|glu ~ .|。

\subsection{先看数据}
任何分析之前先做三件事：\Verb|str()| 看每列的类型，\Verb|summary()| 看量级和范围，\Verb|head()| 看前几行。数值变量被读成了字符（chr）、单位差了 10 倍、出现了不合理的极端值，都能在这一步发现。

\begin{RCode}[R · 第一眼]
str(d); summary(d); head(d)
sapply(d[, -1], mean); sapply(d[, -1], sd)   # 每列的均数、标准差
round(cor(d[, -1]), 3)                        # 相关矩阵
colSums(is.na(d))                             # 每列缺失数
\end{RCode}

取子集的几种写法在诊断时会反复用到：\Verb|d[d$age > 70, ]| 按条件取行，\Verb|subset(d, age > 70, select = c(id, glu))| 同时选列，\Verb|which(d$glu == max(d$glu))| 返回行号，\Verb|d[-149, ]| 删掉第 149 行（删点做敏感性分析就是这样写的）。

\subsection{因子：分组变量与参照组}
分类变量在 R 里要是因子（factor）。建模时，因子的\textbf{第一个水平}就是参照组，回归系数都是“其他水平减参照组”。所以水平顺序必须自己掌握。

\begin{RCode}[R · cut、factor 与 relevel]
d$agegrp <- cut(d$age, breaks = c(29, 45, 58, 79),         # 区间 (29,45]、(45,58]、(58,79]
                labels = c("30-45", "46-58", "59-78"))
table(d$agegrp)                                           # 56 / 94 / 50
levels(d$agegrp)                                          # 第一个就是参照组
d$agegrp2 <- relevel(d$agegrp, ref = "59-78")             # 只换参照组
d$agegrp3 <- factor(d$agegrp, levels = c("46-58", "30-45", "59-78"))  # 指定全部顺序
bp$group  <- factor(bp$group, levels = c("方案A", "方案B"), labels = c("A", "B"))
\end{RCode}

\begin{notebox}[易错：cut() 生成的水平名]
\Verb|cut()| 不给 \Verb|labels| 时，水平名形如 \Verb|"(-Inf,110]"|、\Verb|"(110, Inf]"|——注意后者逗号后面有一个空格（R 对 \Verb|Inf| 的格式化造成的）。第三讲课堂代码 \Verb|factor(..., levels = c("(110, Inf]", "(-Inf,110]"))| 必须一字不差，少了空格就会把整组变成 \Verb|NA|。最稳妥的办法是先 \Verb|levels(x)| 看一眼，再复制粘贴。
\end{notebox}

\subsection{分组汇总}
\begin{RCode}[R · 分组汇总的四种写法]
tapply(d$glu, d$agegrp, mean)                                  # 向量
aggregate(glu ~ agegrp, data = d, FUN = mean)                  # 数据框
aggregate(cbind(base, post) ~ group, data = bp, FUN = mean)    # 多个变量
aggregate(glu ~ agegrp, data = d,
          FUN = function(x) c(mean = mean(x), sd = sd(x), n = length(x)))
sapply(split(d, d$agegrp), function(s) coef(lm(glu ~ bmi, data = s))[2])  # 分层回归
\end{RCode}
最后一行先用 \Verb|split()| 把数据按组拆成列表，再对每一组拟合回归、取斜率，第一讲“分层看斜率”、第三讲“各组单独拟合斜率”、作业四“各组斜率”都是这一行。

\subsection{宽型与长型}
重复测量数据常以宽型录入（一人一行，每个时点一列），而 \Verb|aov_ez()|、\Verb|emmeans()|、\Verb|ggplot()| 都要长型（一人每个时点一行）。转换用 tidyr：

\begin{RCode}[R · pivot\_longer]
library(tidyr)
dm <- read.csv("data/dm.csv", fileEncoding = "UTF-8")
dl <- pivot_longer(dm, cols = c(hba1c_0, hba1c_3, hba1c_6),
                   names_to = "time", values_to = "hba1c")
dl$time <- factor(dl$time, levels = c("hba1c_0", "hba1c_3", "hba1c_6"),
                  labels = c("T0", "M3", "M6"))       # 时间必须是因子，且顺序要对
nrow(dm); nrow(dl)                                    # 40 → 120
pivot_wider(dl, names_from = time, values_from = hba1c)   # 转回宽型
\end{RCode}

\subsection{模型公式}
\Verb|lm()|、\Verb|aov()|、\Verb|t.test()|、\Verb|aggregate()| 都用同一套公式语法，波浪线左边是因变量，右边描述设计矩阵 $X$ 的列：

\begin{table}[htbp]
\centering\small
\caption{常用公式写法}
\begin{tabular}{@{}ll@{}}
\toprule
写法 & 含义\\
\midrule
\Verb|y ~ x1 + x2| & 两个主效应\\
\Verb|y ~ .| & 数据框中除 y 外的全部变量\\
\Verb|y ~ a * b| & 等于 \Verb|a + b + a:b|，主效应加交互\\
\Verb|y ~ a:b| & 只有交互项\\
\Verb|y ~ x - 1| & 去掉截距（分类变量时系数即各组均数）\\
\Verb|y ~ x + I(x^2)| & 公式里做算术要包 \Verb|I()|，否则 \Verb|^| 被理解为交互展开\\
\Verb|log(y) ~ x| & 对因变量变换\\
\Verb|update(fit, . ~ . - x1)| & 在原模型上删变量；\Verb|. ~ . + I(x^2)| 加二次项\\
\Verb|update(fit, data = d[-149, ])| & 换数据重拟合（删点敏感性分析）\\
\Verb|reformulate(c("a","b"), "y")| & 用字符串拼出 \Verb|y ~ a + b|，循环建模时用\\
\bottomrule
\end{tabular}
\end{table}

\Verb|model.matrix(fit)| 可以把公式实际生成的设计矩阵打印出来。第三讲“把研究设计翻译成 X 的列”，在 R 里就是看这个矩阵。

\subsection{从模型对象里取结果}
\Verb|lm()| 的返回值是一个列表，\Verb|summary(fit)| 也是。写报告、做表、手算核对时，用提取函数比从屏幕上抄数字可靠：

\begin{table}[htbp]
\centering\footnotesize
\caption{lm 对象的常用提取函数}
\label{tab:extract}
\begin{tabular}{@{}ll@{\qquad}ll@{}}
\toprule
函数 & 取出什么 & 函数 & 取出什么\\
\midrule
\Verb|coef(fit)| & 系数 $b$ & \Verb|fitted(fit)| & 拟合值 $\hat y$\\
\Verb|confint(fit)| & 系数的 95\% CI & \Verb|resid(fit)| & 残差 $e$\\
\Verb|vcov(fit)| & $s^2(X^\top X)^{-1}$ & \Verb|hatvalues(fit)| & 杠杆值 $h_{ii}$\\
\Verb|sigma(fit)| & 剩余标准差 $s$ & \Verb|rstandard(fit)| & 内学生化残差\\
\Verb|df.residual(fit)| & $n-p-1$ & \Verb|rstudent(fit)| & 外学生化残差\\
\Verb|deviance(fit)| & 残差平方和 SSE & \Verb|cooks.distance(fit)| & Cook 距离\\
\Verb|nobs(fit)| & 样本量 & \Verb|model.matrix(fit)| & 设计矩阵 $X$\\
\Verb|logLik(fit)| & 对数似然 & \Verb|AIC(fit)|、\Verb|BIC(fit)| & 信息准则\\
\Verb|coef(summary(fit))| & 系数表（矩阵） & \Verb|extractAIC(fit)| & step() 用的简略 AIC\\
\Verb|summary(fit)$r.squared| & $R^2$ & \Verb|summary(fit)$adj.r.squared| & 校正 $R^2$\\
\Verb|summary(fit)$fstatistic| & 整体 $F$ 及自由度 & \Verb|anova(fit)| & 方差分析表（I 型）\\
\bottomrule
\end{tabular}
\end{table}

例如 \Verb!coef(summary(fit))["bmi", "Pr(>|t|)"]! 直接取出 bmi 的 P 值；\Verb|coef(summary(fit))[, "t value"]^2| 一次算出各变量的偏 $F$。

\subsection{不用 attach()}
课堂代码常写 \Verb|attach(d)|，之后直接用列名。它的问题是：先后 attach 两份含同名变量的数据，后面的会被前面的遮盖；对 \Verb|d| 做的修改不会同步到已 attach 的副本。第二讲的 \Verb|回归诊断.R| 连续 attach 了三份数据，就是这个隐患。整理版一律改为在函数里写 \Verb|data = d|，临时计算用 \Verb|with(d, cor(bmi, glu))|。

\subsection{保存结果}
\begin{RCode}[R · 存表、存图、存模型]
write.csv(tab, "output/table.csv", row.names = FALSE)
pdf("output/scatter.pdf", width = 5, height = 4)    # 矢量图，插进 LaTeX/Word 都清楚
plot(glu ~ bmi, data = d); dev.off()                 # dev.off() 才真正写完文件
saveRDS(fit, "output/fit.rds"); fit <- readRDS("output/fit.rds")
\end{RCode}

\section{多重线性回归：拟合、读输出与报告}
\srcnote{课件：第一讲第 3–19 页\quad 课堂代码：R上手\_课件一.R、回归2.R\quad 脚本：\texttt{scripts/02\_多重回归\_拟合与解读.R}\quad 数据：datasetA（n = 200，模拟）}

\subsection{粗效应与调整效应}
第一讲的主线是“BMI 对血糖的效应，扣掉年龄之后还剩多少”。在 R 里就是两个模型：

\begin{RCode}[R · 两个模型]
d <- read.csv("data/datasetA.csv")
fit1 <- lm(glu ~ bmi, data = d)          # 粗效应
fit2 <- lm(glu ~ age + bmi, data = d)    # 调整年龄后的偏回归系数
rbind(crude    = coef(summary(fit1))["bmi", ],
      adjusted = coef(summary(fit2))["bmi", ])
\end{RCode}
\begin{ROut}
> rbind(crude = ..., adjusted = ...)
           Estimate Std. Error  t value     Pr(>|t|)
crude    0.10630299 0.01289497 8.243759 2.271622e-14
adjusted 0.07678527 0.01196717 6.416326 1.016509e-09
\end{ROut}
0.106 降到 0.077，粗效应比调整效应高 38\%（\Verb|coef(fit1)["bmi"] / coef(fit2)["bmi"] - 1|）。两个都高度显著，但只有后者是“年龄固定时”的 BMI 效应。

\paragraph{分层的写法。}课件第 5 页把年龄切成三层看斜率，代码是 \Verb|cut()| 加 \Verb|split()|；“三条线能否合成一条共斜率”是一个嵌套模型比较：

\begin{RCode}[R · 分层斜率与平行性检验]
d$agegrp <- cut(d$age, breaks = c(29, 45, 58, 79), labels = c("30-45", "46-58", "59-78"))
sapply(split(d, d$agegrp), function(s) coef(lm(glu ~ bmi, data = s))["bmi"])
m_common <- lm(glu ~ bmi + agegrp, data = d)    # 三层共用一个斜率
m_sep    <- lm(glu ~ bmi * agegrp, data = d)    # 三层各自的斜率
anova(m_common, m_sep)
\end{RCode}
\begin{ROut}
> sapply(...)
 30-45.bmi  46-58.bmi  59-78.bmi
0.10124182 0.07102748 0.10060479
> anova(m_common, m_sep)
  Res.Df    RSS Df Sum of Sq      F Pr(>F)
1    196 64.239
2    194 63.777  2   0.46247 0.7034 0.4962
\end{ROut}
$P = 0.50$，三条斜率的差别在抽样误差之内，\Verb|coef(m_common)["bmi"]| 给出共斜率 0.086。这个"先拟合含交互的大模型，再用 \Verb|anova(小, 大)| 检验交互"的套路，在第 7 章协方差分析的斜率同质性检验里会原样再出现一次。

\subsection{summary() 里该读的数}
\begin{ROut}[summary(fit2) 的输出]
> summary(fit2)
Coefficients:
            Estimate Std. Error t value Pr(>|t|)
(Intercept) 1.654476   0.295334   5.602 7.06e-08 ***
age         0.029963   0.003894   7.695 6.65e-13 ***
bmi         0.076785   0.011967   6.416 1.02e-09 ***

Residual standard error: 0.5393 on 197 degrees of freedom
Multiple R-squared:  0.4276,	Adjusted R-squared:  0.4218
F-statistic: 73.58 on 2 and 197 DF,  p-value: < 2.2e-16
\end{ROut}

\begin{table}[htbp]
\centering\small
\caption{summary() 输出的七个数及提取方式}
\begin{tabular}{@{}lll@{}}
\toprule
输出位置 & 含义 & 取出来的写法\\
\midrule
Estimate & 偏回归系数 $b$ & \Verb|coef(fit2)|\\
Std. Error & $b$ 的标准误 & \Verb|sqrt(diag(vcov(fit2)))|\\
t value、Pr(>|t|) & $t=b/\mathrm{SE}$ 及其双侧 P & \Verb|coef(summary(fit2))[, 3:4]|\\
Residual standard error & 剩余标准差 $s$，自由度 $n-p-1$ & \Verb|sigma(fit2)|、\Verb|df.residual(fit2)|\\
Multiple R-squared & $R^2$ & \Verb|summary(fit2)$r.squared|\\
Adjusted R-squared & 校正 $R^2$ & \Verb|summary(fit2)$adj.r.squared|\\
F-statistic & 整个方程的 $F$ 检验 & \Verb|summary(fit2)$fstatistic|\\
\bottomrule
\end{tabular}
\end{table}

整体检验的 P 值 summary 只显示“< 2.2e-16”，要精确值可以 \Verb|pf(F, df1, df2, lower.tail = FALSE)|，本例为 $1.36\times10^{-24}$。各项都可以手算核对：

\begin{RCode}[R · 手算核对]
b <- coef(fit2); se <- sqrt(diag(vcov(fit2)))
cbind(t = b / se, P = 2 * pt(-abs(b / se), df.residual(fit2)))   # 与 summary 相同
n <- 200; p <- 2
1 - (1 - summary(fit2)$r.squared) * (n - 1) / (n - p - 1)        # 校正 R² = 0.4218
deviance(fit2) / df.residual(fit2); sigma(fit2)^2                 # MSE = 0.2908
\end{RCode}

\subsection{方差分析表：顺序很重要}
\begin{ROut}
> anova(fit2)                          # age 先进
           Df Sum Sq Mean Sq F value    Pr(>F)
age         1 30.824 30.8237 105.987 < 2.2e-16 ***
bmi         1 11.973 11.9730  41.169 1.017e-09 ***
Residuals 197 57.292  0.2908
> anova(lm(glu ~ bmi + age, data = d)) # bmi 先进
bmi         1 25.575 25.5754  87.941 < 2.2e-16 ***
age         1 17.221 17.2213  59.215 6.646e-13 ***
Residuals 197 57.292  0.2908
\end{ROut}
\Verb|anova(fit)| 给的是 I 型（顺序）平方和：每一行是“在它前面的变量都已入模后”新增的解释量。bmi 写在后面时 SS = 11.97，是“调整年龄后”的贡献；写在前面时 SS = 25.58，是它单独的贡献。残差行永远不变。最后一个变量那一行的 $F$（41.17）等于它在 summary 里的 $t^2$（$6.416^2$）。这一点在协方差分析里会变成一个真实的陷阱（第 \ref{sec:ancova-order} 节）。

\subsection{系数的置信区间}
\begin{RCode}[R · confint 与手算]
confint(fit2)                        # 默认 95%；level = 0.99 改水平
tq <- qt(0.975, df.residual(fit2))   # 1.972
cbind(b - tq * se, b + tq * se)      # bmi：0.053 — 0.100
\end{RCode}

\subsection{预测：均数的 CI 与个体的 PI}
\begin{RCode}[R · predict]
new <- data.frame(age = 52, bmi = 24)     # 列名必须与模型中完全一致
predict(fit2, newdata = new, interval = "confidence")
predict(fit2, newdata = new, interval = "prediction")
pr <- predict(fit2, newdata = new, se.fit = TRUE)
c(se_mean = pr$se.fit, se_indiv = sqrt(sigma(fit2)^2 + pr$se.fit^2))
\end{RCode}
\begin{ROut}
> predict(fit2, newdata = new, interval = "confidence")
       fit      lwr      upr
1 5.055392 4.980162 5.130622
> predict(fit2, newdata = new, interval = "prediction")
1 5.055392 3.989228 6.121556
> c(se_mean = ..., se_indiv = ...)
   se_mean   se_indiv
0.03814755 0.54062936
\end{ROut}
两个区间中心相同，宽度差 14 倍，差别全在标准误：均数的 CI 只含估计误差（0.038），个体的 PI 还要加上个体本身的波动 $s$（$\sqrt{0.539^2+0.038^2}=0.541$）。\Verb|newdata| 可以一次放多行，得到多个点的预测。

\begin{notebox}[易错：newdata 的列名]
报错 \Verb|'newdata' had 1 row but variables found have 200 rows| 几乎都是 \Verb|newdata| 的列名和模型里的变量名不一致（或少了某个自变量），于是 R 退回去用原数据。课堂代码 \Verb|筛选自变量.R| 里先把变量改名为 \Verb|x2, x3, x4| 再建模，\Verb|newdata| 也必须用 \Verb|x2, x3, x4|。建模时始终写 \Verb|data =|、不要 \Verb|attach()|，这个问题就少很多。
\end{notebox}

\subsection{标准化偏回归系数}
\begin{RCode}[R · 三种写法]
b[-1] * sapply(d[, c("age", "bmi")], sd) / sd(d$glu)    # b* = b × SD(x)/SD(y)
dz <- as.data.frame(scale(d[, c("age", "bmi", "glu")]))
coef(lm(glu ~ age + bmi, data = dz))                    # y、x 都标准化后回归
coef(lm(glu ~ scale(age) + scale(bmi), data = d))       # 错：只标准化了 x
\end{RCode}
前两种都得到 age 0.438、bmi 0.365（原始系数是 bmi 大，标准化后 age 大，排序反转）；第三种得到 0.311、0.259，是只除了 SD(x) 的“半标准化”系数，差一个 SD(y) 倍。课堂代码 \Verb|R上手_课件一.R| 专门提醒过这一点。标准化系数不改变任何检验结论，作业一要求的 95\% CI 也按同一比例缩放：\Verb|confint(fit)[-1, ] * sx / sy|。

\subsection{三种相关：r、复相关 R、偏相关}
\begin{RCode}[R · 相关]
cor(d$glu, d$bmi)                                   # 简单相关 0.5055
cor(d$glu, fitted(fit2)); sqrt(summary(fit2)$r.squared)   # 复相关 R = 0.6539
e_y <- resid(lm(glu ~ age, data = d))               # 偏相关：两组残差的相关
e_x <- resid(lm(bmi ~ age, data = d))
cor(e_y, e_x)                                       # 0.4158
library(ppcor)
pcor.test(d$glu, d$bmi, d$age)                      # estimate 0.4158，statistic 6.416
\end{RCode}
\Verb|pcor.test(x, y, z)| 的 statistic 就是偏相关的 $t$，它与全模型中 bmi 的 $t$（6.416）完全相同；也可以手算 $r\sqrt{(n-2-q)/(1-r^2)}$，$q$ 为控制变量个数。“残差对残差”的写法要记住，第 4 章的偏回归图、作业一思考题的偏相关都是它。

\subsection{lm() 背后的矩阵运算}
\begin{RCode}[R · 正规方程、帽子矩阵与系数协方差阵]
X <- model.matrix(fit2); y <- d$glu
XtX_inv <- solve(t(X) %*% X)
bhat <- XtX_inv %*% t(X) %*% y            # b = (X'X)^{-1} X'y，与 coef(fit2) 相同
H <- X %*% XtX_inv %*% t(X)               # 帽子矩阵，yhat = H y
h <- diag(H)                              # 杠杆值，等于 hatvalues(fit2)
mean(h)                                   # = (p + 1)/n = 0.015
which.max(h)                              # 51 号：x 方向离中心最远
s2 <- sum((y - X %*% bhat)^2) / (200 - 3)
s2 * XtX_inv                              # 等于 vcov(fit2)
\end{RCode}
\Verb|%*%| 是矩阵乘法，\Verb|t()| 转置，\Verb|solve()| 求逆，\Verb|diag()| 取对角线。\Verb|crossprod(X)| 是 \Verb|t(X) %*% X| 的快捷写法。

\subsection{对数似然、AIC 与 BIC}
\begin{RCode}[R · 手算与 R 函数对照]
sse <- deviance(fit2); n <- 200; K <- 4          # K = 截距 + 2 个系数 + sigma²
ll <- -n/2 * log(2*pi) - n/2 * log(sse/n) - n/2
c(ll, logLik(fit2))                              # -158.77
c(-2*ll + 2*K, AIC(fit2))                        # 325.55
c(-2*ll + log(n)*K, BIC(fit2))                   # 338.74
extractAIC(fit2)                                 # 3, -244.03：step() 用的简略版
n * log(sse/n) + 2 * 3                           # 简略 AIC = n·ln(SSE/n) + 2(p+1)
AIC(fit2) - extractAIC(fit2)[2]                  # 569.58 = n(1 + ln 2pi) + 2
\end{RCode}
\Verb|AIC()| 与 \Verb|extractAIC()| 只差一个对同一数据恒定的常数，比较模型时结论相同，但数值不能混着报。\Verb|step()| 打印的 “AIC=” 是 \Verb|extractAIC()|。

\subsection{整理成论文表格}
\begin{RCode}[R · 结果表]
tab <- data.frame(term = names(coef(fit2))[-1], b = coef(fit2)[-1],
                  lo = confint(fit2)[-1, 1], hi = confint(fit2)[-1, 2],
                  P = coef(summary(fit2))[-1, 4])
tab$report <- sprintf("%.3f (%.3f, %.3f)", tab$b, tab$lo, tab$hi)
tab$P <- format.pval(tab$P, digits = 3, eps = 0.001)
tab[, c("term", "report", "P")]
\end{RCode}
\begin{ROut}
> tab[, c("term", "report", "P")]
    term               report      P
age  age 0.030 (0.022, 0.038) <0.001
bmi  bmi 0.077 (0.053, 0.100) <0.001
\end{ROut}

\begin{resultbox}[写法]
在同时调整年龄后，BMI 每增加 1 kg/m²，空腹血糖平均升高 0.077 mmol/L（95\% CI 0.053–0.100，$P<0.001$）；模型 $R^2=0.428$，剩余标准差 0.539 mmol/L，$n=200$。报告时写清调整了哪些变量，给效应量和区间，不要只给 P 值，也不要把标准化系数当作带单位的效应来解释。
\end{resultbox}

\section{嵌套模型检验与变量筛选}
\srcnote{课件：第一讲第 20–24 页\quad 课堂代码：筛选自变量.R、stepwise.R、回归诊断.R 前半\quad 脚本：\texttt{scripts/03\_假设检验与变量筛选.R}\quad 数据：datasetC（n = 200，5 个候选变量）、diabetes27\_hw（n = 27）}

\subsection{增量 F 检验：anova(小模型, 大模型)}
所有“去掉某些变量行不行”的问题，都是比较两个嵌套模型。R 的通用写法是 \Verb|anova(小, 大)|：

\begin{RCode}[R · 偏 F 检验]
dd   <- read.csv("data/diabetes27_hw.csv")[, -1]
full <- lm(glucose ~ tc + tg + insulin + hba1c, data = dd)
m3   <- lm(glucose ~ tg + insulin + hba1c, data = dd)
anova(m3, full)                                       # 检验 tc
F_tc <- (deviance(m3) - deviance(full)) / 1 / (deviance(full) / df.residual(full))
pf(F_tc, 1, df.residual(full), lower.tail = FALSE)    # 手算，同上
coef(summary(full))[, "t value"]^2                    # 单个变量：偏 F = t²
anova(lm(glucose ~ insulin + hba1c, data = dd), full) # 同时检验 tc 与 tg
\end{RCode}
\begin{ROut}
> anova(m3, full)
Model 1: glucose ~ tg + insulin + hba1c
Model 2: glucose ~ tc + tg + insulin + hba1c
  Res.Df    RSS Df Sum of Sq      F Pr(>F)
1     23 89.454
2     22 88.841  1   0.61285 0.1518 0.7006
\end{ROut}
输出里 RSS 一列就是两个模型的残差平方和，$F=(89.454-88.841)/1\div(88.841/22)=0.152$。作业一要求“写出检验过程，包括 F 值、自由度、P 值”，三样都在这张表里：$F(1,22)=0.152$，$P=0.70$。

\subsection{drop1() 与 add1()}
\Verb|drop1(fit, test = "F")| 对模型里每个变量各做一次“删掉它”的偏 F 检验，\Verb|add1(fit, scope, test = "F")| 对候选变量各做一次“加上它”的检验。后退法的每一步看 \Verb|drop1|，前进法的每一步看 \Verb|add1|。不写 \Verb|test = "F"| 时只给 AIC，不给 P 值（课堂代码 \Verb|stepwise.R| 后半就是这种情况，所以要手算 F）。

\subsection{I 型、II 型、III 型平方和}
\begin{RCode}[R · 顺序与平方和类型]
anova(lm(glucose ~ tc + tg, data = dd))     # I 型：tc 的 SS 是它单独的贡献
anova(lm(glucose ~ tg + tc, data = dd))     # 顺序一换，两行都变
cor(dd$tc, dd$tg)                           # 0.63：两者相关才会这样
library(car)
Anova(lm(glucose ~ tc + tg, data = dd), type = "II")   # 与顺序无关
\end{RCode}
\begin{table}[htbp]
\centering\small
\caption{三种平方和的区别与 R 写法}
\begin{tabular}{@{}p{2.2cm}p{7.2cm}p{4.6cm}@{}}
\toprule
类型 & 每一行检验什么 & R\\
\midrule
I 型（顺序） & 在排在它前面的项之后新增的解释量 & \Verb|anova(fit)|\\
II 型 & 在其他所有\textbf{同级或更低级}项之后的贡献（不调整含它的交互项） & \Verb|car::Anova(fit, type = "II")|\\
III 型 & 在其他所有项（含交互）之后的贡献 & 先设和为零编码，再 \Verb|Anova(fit, type = "III")|\\
\bottomrule
\end{tabular}
\TableNote{无交互项时，II 型与 III 型相同，都等于 \texttt{drop1()} 的偏 F。平衡的析因设计中三者相同；不平衡时三者不同（第 6.9 节）。}
\end{table}

\subsection{为什么不能用 R² 选模型}
\begin{RCode}[R · 一组嵌套模型的五个指标]
dc <- read.csv("data/datasetC.csv")[, -1]
forms <- list("1" = glu ~ 1, "age" = glu ~ age, "age+bmi" = glu ~ age + bmi,
              "age+bmi+exercise" = glu ~ age + bmi + exercise,
              "+sleep" = glu ~ age + bmi + exercise + sleep,
              "full" = glu ~ age + bmi + exercise + sleep + waist)
fits <- lapply(forms, lm, data = dc)
t(sapply(fits, function(m) c(R2 = summary(m)$r.squared,
  adjR2 = summary(m)$adj.r.squared, AIC = AIC(m), BIC = BIC(m))))
\end{RCode}
\begin{ROut}
> t(sapply(fits, ...))
                     R2  adjR2      AIC      BIC
1                0.0000 0.0000 428.6045 435.2011
age              0.2791 0.2755 365.1542 375.0492
age+bmi          0.4191 0.4132 323.9729 337.1662
age+bmi+exercise 0.4692 0.4611 307.9146 324.4062
+sleep           0.4702 0.4594 309.5350 329.3249
full             0.4712 0.4576 311.1739 334.2621
\end{ROut}
$R^2$ 一路上升；校正 $R^2$、AIC、BIC 都在 age+bmi+exercise 处最优，与课件第 23 页的表一致。\Verb|lapply(列表, lm, data = dc)| 一次拟合一组模型，是比较多个模型时很省事的写法。

\subsection{全子集：leaps::regsubsets()}
\begin{RCode}[R · 全子集回归]
library(leaps)
all_sub <- regsubsets(glu ~ ., data = dc, nbest = 1)   # 每个变量个数下保留最好的 1 个
ss <- summary(all_sub)
ss$which                                     # 每行入选了哪些变量（TRUE/FALSE 矩阵）
cbind(adjR2 = ss$adjr2, Cp = ss$cp, BIC = ss$bic)
which.max(ss$adjr2); which.min(ss$cp); which.min(ss$bic)   # 都是第 3 行
coef(all_sub, 3)                             # 第 3 行模型的系数
plot(all_sub, scale = "adjr2")               # 黑格 = 入选
\end{RCode}
注意 \Verb|summary(regsubsets(...))| 返回的行号是模型的序号，不是变量个数（\Verb|nbest = 5| 时第 11 行是“3 个变量里排第 1”的模型）；课件笔记例 15-1 里 \Verb|which.max(...$adjr2)| 返回 11 就是这个意思。$C_p$ 的理想值是 $p+1$，本例 3 个变量时 $C_p=2.72$。\Verb|regsubsets| 返回的不是 lm 对象，不能直接 \Verb|vif()|、\Verb|summary()| 出系数表；选好后要用选中的变量重新 \Verb|lm()| 一次。

\subsection{按 AIC 的逐步法：step()}
\begin{RCode}[R · step 的四种用法]
m0    <- lm(glu ~ 1, data = dc)
mfull <- lm(glu ~ age + bmi + waist + sleep + exercise, data = dc)
step(mfull, direction = "backward")                                  # 后退
step(m0, scope = list(lower = ~ 1, upper = formula(mfull)), direction = "forward")
step(m0, scope = list(lower = ~ 1, upper = formula(mfull)), direction = "both")
step(m0, scope = list(lower = ~ 1, upper = formula(mfull)), direction = "both",
     k = log(nrow(dc)))                                              # 按 BIC
step(mfull, scope = list(lower = ~ age, upper = formula(mfull)))     # 强制保留 age
\end{RCode}
\begin{ROut}
> step(m0, scope = ..., direction = "forward")
Start:  AIC=-140.97
glu ~ 1
           Df Sum of Sq    RSS     AIC
+ age       1    27.310 70.542 -204.42
+ bmi       1    22.137 75.715 -190.27
+ waist     1    16.882 80.970 -176.85
+ exercise  1    10.647 87.206 -162.01
<none>                  97.853 -140.97
+ sleep     1     0.034 97.819 -139.04
\end{ROut}
读法：每一步列出“做这个动作后的 AIC”，\Verb|<none>| 是不动，取 AIC 最小的那一行，直到 \Verb|<none>| 最小为止。四种方式在本例都停在 age+bmi+exercise。几点要点：

\begin{itemize}
\item \Verb|scope| 的 \Verb|lower| 写进去的变量不会被删，混杂因素应当这样强制保留，而不是交给 AIC 决定。
\item \Verb|k = log(n)| 把惩罚从 2 换成 $\ln n$，就是按 BIC 选。
\item 这里的 AIC 是 \Verb|extractAIC()|。课堂代码 \Verb|回归诊断.R| 里那几行 \Verb|-140.97 + 200*(1+log(2*pi)) + 2| 就是把它换算成 \Verb|AIC()| 的数值：$-140.97+200(1+\ln2\pi)+2=428.60$，正是上表空模型的 AIC。
\end{itemize}

\subsection{按 P 值逐步：手工一步步走}
课件和作业用的是按 P 值的前进、后退、逐步法（$\alpha_{\text{入}}$、$\alpha_{\text{出}}$），R 的 \Verb|step()| 不做这个，可以用 \Verb|add1/drop1| 手工走，或者用 olsrr 自动走。

\begin{RCode}[R · 前进法与后退法（diabetes27）]
sc <- ~ tc + tg + insulin + hba1c
s0 <- lm(glucose ~ 1, data = dd)
add1(s0, sc, test = "F")                    # 第 1 步：hba1c 的 P 最小，进入
s1 <- update(s0, . ~ . + hba1c)
add1(s1, sc, test = "F")                    # 第 2 步：tc 进入（P = 0.03）
s2 <- update(s1, . ~ . + tc)
drop1(s2, test = "F")                       # 逐步法：进完一步要回头看能否删
add1(s2, sc, test = "F")                    # 第 3 步：最小 P = 0.088 > 0.05，停
drop1(full, test = "F")                     # 后退法：tc 的 P 最大（0.70），删
drop1(update(full, . ~ . - tc), test = "F") # 余下三个都 < 0.05，停
\end{RCode}
前进法（$\alpha=0.05$）停在 tc + hba1c，后退法停在 tg + insulin + hba1c——起点不同，终点就不同。后退法的循环可以写成一个函数，脚本 §3.7 的 \Verb|backward_p()| 就是作业一 bootstrap 用的那个。

\begin{RCode}[R · olsrr 自动逐步]
library(olsrr)
ols_step_forward_p(full, p_val = 0.05)         # 前进：p_val
ols_step_backward_p(full, p_val = 0.05)        # 后退：p_val
k <- ols_step_both_p(full, p_enter = 0.10, p_remove = 0.15)
k$model                                     # 最终模型（lm 对象，可以接着 summary）
\end{RCode}
olsrr（本机 0.7.1，课堂代码的写法也是如此）里前进、后退两个函数的阈值参数叫 \Verb|p_val|，逐步法才是 \Verb|p_enter|、\Verb|p_remove|。参数名写错不会报错，只会被 \Verb|...| 吞掉、悄悄用默认值 0.3，所以跑完要看一眼输出里的入选标准。

\subsection{筛选结果稳不稳：bootstrap}
\begin{RCode}[R · 重抽样 500 次看后退法的结果]
set.seed(2026)
sel <- replicate(500, {
  i <- sample(nrow(dd), replace = TRUE)
  paste(backward_p(dd[i, ], c("tc", "tg", "insulin", "hba1c"), "glucose"), collapse = "+")
})
head(sort(table(sel), decreasing = TRUE) / 500)
sapply(c("tc", "tg", "insulin", "hba1c"), function(v) mean(grepl(v, sel)))
\end{RCode}
\begin{ROut}
> head(sort(table(sel), decreasing = TRUE) / 500)
   tg+insulin+hba1c          tg+insulin            tc+hba1c               hba1c
              0.286               0.162               0.128               0.120
> sapply(...)           # 各变量被选入的比例
     tc      tg insulin   hba1c
  0.262   0.544   0.652   0.712
\end{ROut}
全样本选出的模型在重抽样中只出现 29\%。27 例、4 个相关的自变量，筛选结果本身就很不稳定，这正是课件说“筛过的 P 值要打折”的原因。\Verb|replicate(B, {...})| 加 \Verb|sample(n, replace = TRUE)| 是 bootstrap 最简单的写法。

\begin{notebox}[筛选的几条规矩]
先按专业知识定下必须调整的变量（放进 \Verb|scope| 的 \Verb|lower|），再让准则处理“可加可不加”的变量；写清起点、方向和准则；筛选后报告的 P 值和 CI 偏乐观，要在方法部分交代“变量经逐步法筛选”；样本小、自变量相关时，用 bootstrap 看一眼结果是否稳定。
\end{notebox}

\section{回归诊断}
\srcnote{课件：第二讲第 3–9 页、第 16–21 页\quad 脚本：\texttt{scripts/04\_回归诊断.R}\quad 数据：datasetB（n = 150，模拟数据，埋了异方差、共线性和 149、150 两个异常个体）}

第二讲的例子是一个“看起来很好”的模型：$R^2=0.459$，$F=41.35$，$P<0.001$。诊断要回答的是：这些数字能不能信。

\begin{RCode}[R · 模型]
library(car); library(MASS); library(lmtest); library(sandwich)
d <- read.csv("data/datasetB.csv")          # 注意 id 在最后一列
m <- lm(glu ~ age + bmi + waist, data = d)
n <- nrow(d); p <- 3                        # p = 自变量个数
\end{RCode}

\begin{table}[htbp]
\centering\small
\caption{诊断清单：每一项用什么函数、看什么、本例结果}
\label{tab:diag}
\begin{tabular}{@{}p{1.7cm}p{4.6cm}p{3.6cm}p{4.4cm}@{}}
\toprule
条件 & R & 看什么 & 本例\\
\midrule
线性 & \Verb|plot(m, which = 1)|、\Verb|residualPlots(m)|、\Verb|crPlots(m)| & 平滑线是否贴着 0；曲率检验 P & 大体可以；bmi、waist 曲率 P≈0.01–0.04\\
独立 & 研究设计；有序数据才用 \Verb|durbinWatsonTest(m)| & 抽样是否独立 & 横断面数据，DW 无意义\\
正态 & \Verb|plot(m, which = 2)|、\Verb|shapiro.test(resid(m))| & Q-Q 图两端 & W = 0.900，P = 1.4e-8，重尾\\
等方差 & \Verb|plot(m, which = 3)|、\Verb|ncvTest(m)| & 点带是否喇叭形 & χ² = 57.6，P = 3.2e-14\\
共线性 & \Verb|vif(m)|、条件数 & VIF > 10 & bmi 16.0、waist 16.1（第 5 章）\\
强影响点 & \Verb|hatvalues|、\Verb|rstudent|、\Verb|cooks.distance|、\Verb|influencePlot| & 杠杆 × 残差 & 149 号 D = 0.86；150 号 t = 6.85\\
\bottomrule
\end{tabular}
\end{table}

\subsection{四张诊断图}
\Verb|plot(m)| 依次输出四张图（\Verb|par(mfrow = c(2, 2))| 放在一页）：

\begin{figure}[htbp]
\centering
\includegraphics[width=0.92\textwidth]{figures/fig1_diag4.pdf}
\caption{datasetB 全模型的四张诊断图。左上：残差–拟合值；右上：正态 Q-Q；左下：尺度–位置；右下：残差–杠杆（虚线为 Cook 距离等高线）}
\label{fig:diag4}
\end{figure}

\begin{itemize}
\item \textbf{which = 1 残差–拟合值}：平滑线是否贴着 0 看线性，点带宽度是否一致看等方差。本例平滑线大体水平，点带左窄右宽。
\item \textbf{which = 2 Q-Q 图}：只看两端。本例右上 150、149 号明显上翘，左下 102 号下坠。
\item \textbf{which = 3 尺度–位置}：纵轴是 $\sqrt{|\text{标准化残差}|}$，平滑线上斜就是方差随拟合值增大。
\item \textbf{which = 5 残差–杠杆}：右上角越过 Cook 等高线的点是强影响点，本例是 149 号。
\end{itemize}
\Verb|plot(m, which = 4)| 是 Cook 距离针状图，\Verb|plot(m, which = 1, id.n = 3)| 可以改标注点的个数。

\subsection{正态性与等方差}
\begin{RCode}[R · 正态与方差齐性]
shapiro.test(resid(m))                   # W = 0.900，P = 1.4e-8
qqPlot(m, id = list(n = 2))              # car 版 Q-Q 图，带置信带，标出 102、150
ncvTest(m)                               # Chisquare = 57.6，P = 3.2e-14
q <- cut(fitted(m), quantile(fitted(m), 0:4 / 4), include.lowest = TRUE)
round(tapply(resid(m), q, sd), 2)        # 按拟合值四等分后各段残差的 SD
\end{RCode}
\begin{ROut}
> round(tapply(resid(m), q, sd), 2)
   [4,4.55] (4.55,4.85] (4.85,5.37] (5.37,6.01]
       0.20        0.34        0.50        0.88
\end{ROut}
最高段残差标准差是最低段的 4 倍多。课件第 7 页表中写的是 0.19 → 0.98，分段方式略有不同，结论一致。\Verb|ncvTest()| 的原假设是“方差齐”；它和 Shapiro 检验一样对样本量很敏感，图和检验要一起看。

\subsection{独立性}
\begin{RCode}[R · Durbin–Watson]
durbinWatsonTest(m)                                   # DW = 1.653，P ~ 0.04（bootstrap，每次略不同）
dwtest(m)                                             # lmtest 版本，单侧
durbinWatsonTest(update(m, data = d[sample(n), ]))    # 打乱行序，结果就变了
\end{RCode}
DW 检验的前提是观测有天然顺序（时间、空间）。本例是横断面体检数据，行序只是录入先后，打乱行序 DW 就变，说明它在这里没有意义。独立性靠研究设计判断：重复测量、整群抽样、同一家庭或医院的人，误差都不独立，要用重复测量方差分析、混合效应模型或 GEE（第 8 章）。

\subsection{线性}
\begin{RCode}[R · 线性的几种查法]
residualPlots(m)        # 每个自变量一张残差图 + 曲率检验；Tukey test 针对整体
crPlots(m)              # 成分 + 残差图：看某个自变量是否需要非线性项
avPlots(m)              # 偏回归图：斜率等于该变量在全模型里的系数
anova(m, update(m, . ~ . + I(age^2)))     # 怀疑谁弯曲，就加谁的二次项检验
\end{RCode}
\begin{ROut}
> residualPlots(m)
           Test stat Pr(>|Test stat|)
age           1.8721          0.06320 .
bmi           2.0911          0.03826 *
waist         2.5092          0.01320 *
Tukey test    1.1185          0.26337
\end{ROut}
\Verb|residualPlots()| 每一行是把该变量的二次项加进模型后的 $t$ 检验，最后一行 Tukey 检验看拟合值的二次项。作业二第一题就是靠这一步发现 hba1c 的二次项显著（$P=0.016$）。偏回归图也可以手工画，原理就是第 2 章“残差对残差”：

\begin{RCode}[R · 手工偏回归图（以 bmi 为例）]
e_y <- resid(lm(glu ~ age + waist, data = d))   # y 中扣掉其他自变量的部分
e_x <- resid(lm(bmi ~ age + waist, data = d))   # bmi 中扣掉其他自变量的部分
plot(e_x, e_y); abline(lm(e_y ~ e_x))
coef(lm(e_y ~ e_x))[2]                          # 0.0671，等于 coef(m)["bmi"]
\end{RCode}

\subsection{方差不齐的三种补救}
\begin{RCode}[R · WLS、稳健标准误、对数变换]
w  <- 1 / fitted(lm(abs(resid(m)) ~ fitted(m)))^2     # 估计每人的 SD，权重 = 1/SD²
mw <- lm(glu ~ age + bmi + waist, data = d, weights = w)
ncvTest(mw)                                             # P = 0.28：补救后再查一遍
coeftest(m, vcov = vcovHC(m, type = "HC3"))             # 系数不变，只换标准误
ml <- lm(log(glu) ~ age + bmi + waist, data = d)
ncvTest(ml)                                             # chi2 = 40.9，改善但未根治
(exp(coef(ml)[-1]) - 1) * 100                           # log(y) 模型：每增 1 单位变化的百分数
\end{RCode}

\begin{table}[htbp]
\centering\small
\caption{三种处理下 bmi 的系数与标准误}
\begin{tabular}{@{}lS[table-format=1.4]S[table-format=1.4]S[table-format=1.3]l@{}}
\toprule
做法 & {系数} & {标准误} & {P} & 说明\\
\midrule
普通 OLS & 0.0671 & 0.0552 & 0.226 & 标准误按方差齐算，不可信\\
加权最小二乘 & 0.0642 & 0.0238 & 0.008 & 对症：系数略变，标准误缩小\\
OLS + HC3 稳健标准误 & 0.0671 & 0.0705 & 0.343 & 系数不动，标准误更保守\\
\bottomrule
\end{tabular}
\TableNote{数据为 datasetB，模型 glu ~ age + bmi + waist。WLS 权重按课件第 21 页的做法估计。}
\end{table}
加权之后 bmi 的 P 值跨过了 0.05，但 waist 依然不显著。这说明 bmi、waist 不显著的主因不在方差，而在共线性（第 5 章）。

\subsection{杠杆值、学生化残差与 Cook 距离}
\label{sec:influence}
\begin{RCode}[R · 影响度量]
inf <- data.frame(id = d$id, h = hatvalues(m), r_int = rstandard(m),
                  r_ext = rstudent(m), D = cooks.distance(m))
head(inf[order(-inf$D), ], 3)                    # 按 Cook 距离排序
c(mean_h = (p + 1) / n, h_cut = 2 * (p + 1) / n, D_cut = 4 / (n - p - 1))
which(hatvalues(m) > 2 * (p + 1) / n)            # 高杠杆
which(abs(rstudent(m)) > 3)                      # 响应方向离群
which(cooks.distance(m) > 4 / (n - p - 1))       # 强影响
influencePlot(m)                                 # 横轴杠杆，纵轴学生化残差，圈大小 = D
outlierTest(m)                                   # 最大 |t| 的 Bonferroni 校正 P
summary(influence.measures(m))                   # DFBETAS、DFFIT、COVRATIO 一次给出
\end{RCode}
\begin{ROut}
> head(inf[order(-inf$D), ], 3)
     id          h     r_int     r_ext          D
149 149 0.29653465  2.863906  2.937793 0.86435133
150 150 0.03165966  5.972340  6.846671 0.29154614
102 102 0.03030998 -3.150737 -3.252448 0.07757414
> c(mean_h = ..., h_cut = ..., D_cut = ...)
    mean_h      h_cut      D_cut
0.02666667 0.05333333 0.02739726
\end{ROut}
149 号（24 岁、BMI 36.2、腰围 119）杠杆值 0.297，是平均值的 11 倍，Cook 距离最大；150 号（70 岁、偏瘦、血糖 8.8）学生化残差 6.85，但待在自变量空间中心，“拽不动”回归线，D 反而较小。残差最大的点不一定是影响最大的点。

\begin{notebox}[几个界值的写法要和自己的 p 对上]
本笔记和课件中 $p$ 指自变量个数，含截距的参数个数是 $p+1$。平均杠杆 $=(p+1)/n$，高杠杆界 $2(p+1)/n$，Cook 距离界 $4/(n-p-1)$。如果写成 \Verb|k <- length(coef(m))|（已含截距），界值就是 \Verb|2*k/n| 和 \Verb|4/(n-k)|，数值相同。这些都只是“筛出来逐个看原始记录”的筛子。
\end{notebox}

三个量都可以手算，公式与课件笔记第 3 章一致：

\begin{RCode}[R · 手算核对]
h <- hatvalues(m); e <- resid(m); s <- sigma(m); k <- p + 1
r  <- e / (s * sqrt(1 - h))                       # 内学生化残差 = rstandard(m)
tt <- r * sqrt((n - k - 1) / (n - k - r^2))       # 外学生化残差 = rstudent(m)
D  <- r^2 / k * h / (1 - h)                       # Cook 距离 = cooks.distance(m)
yhat_del <- predict(update(m, data = d[-149, ]), newdata = d)
sum((fitted(m) - yhat_del)^2) / (k * s^2)         # 定义：删掉 149 号后全部拟合值的移动
\end{RCode}

\subsection{敏感性分析与稳健回归}
\begin{RCode}[R · 删点重拟合与 rlm]
m149 <- update(m, data = d[-149, ]); m150 <- update(m, data = d[-150, ])
sapply(list(all = m, del149 = m149, del150 = m150), function(x)
  c(coef(x), R2 = summary(x)$r.squared, sigma = sigma(x)))
mr  <- rlm(glu ~ age + bmi + waist, data = d)               # Huber M 估计
round(mr$w[c(149, 150)], 3)                                 # 0.375、0.161：异常点被降权
mmm <- rlm(glu ~ age + bmi + waist, data = d, method = "MM")
\end{RCode}
\begin{ROut}
> sapply(...)
                all  del149  del150
(Intercept)  2.4617  3.4314  2.4427
age          0.0314  0.0336  0.0287
bmi          0.0671  0.1241  0.0739
waist       -0.0063 -0.0369 -0.0068
R2           0.4594  0.4791  0.4963
sigma        0.5438  0.5301  0.4744
\end{ROut}
删掉 149 号，bmi 系数从 0.067 跳到 0.124；删掉 150 号，系数几乎不变但残差标准差明显下降。这张“删除前 / 删除后”对照表就是作业二第一题第 4 问要的东西。\Verb|update(m, data = d[-k, ])| 是最简洁的写法；按 id 删更不容易出错：\Verb|d[d$id != 149, ]|。

\begin{resultbox}[强影响点的处理顺序]
先查原始记录（录入错误、单位错误能核实就核实）；确认是真实数据就保留，把含与不含该点的结果都报告；也可以用 \Verb|MASS::rlm()| 做稳健回归作为补充。不能因为删掉之后 P 值更好看就删点。
\end{resultbox}

\subsection{一份可复用的诊断函数}
脚本 §4.11 把上面的检查写成了 \Verb|diag_report(fit)|，拿到新模型可以直接调用：

\begin{ROut}
> diag_report(lm(glucose ~ tc + tg + insulin + hba1c, data = read.csv("data/diabetes27_hw.csv")))
Shapiro-Wilk P = 0.823
ncvTest P      = 0.984
VIF:
     tc      tg insulin   hba1c
   2.19    1.78    1.28    1.27
h > 2k/n      : 7 25
|rstudent| > 3:
Cook D > 4/(n-k): 6 13 25 26
\end{ROut}
这是作业二第一题的全模型：正态和等方差都没问题，共线性不严重，25 号同时是高杠杆点和 Cook 距离最大的点（D = 0.99），删掉它后 tc 的系数由正变负（第 9 章）。

\section{共线性诊断与岭回归}
\srcnote{课件：第二讲第 10–15 页\quad 课堂代码：回归诊断.R 末尾\quad 脚本：\texttt{scripts/05\_共线性与岭回归.R}\quad 数据：datasetB}

\subsection{后果：五个模型并排}
\begin{RCode}[R · 同一批人、不同自变量组合]
x <- c("age", "bmi", "waist")
round(cor(d[, c(x, "glu")]), 3)                  # r(bmi, waist) = 0.968
mods <- list(bmi = lm(glu ~ bmi, d), waist = lm(glu ~ waist, d),
             "age+bmi" = lm(glu ~ age + bmi, d), "age+waist" = lm(glu ~ age + waist, d),
             "age+bmi+waist" = m)
get <- function(fit, v) {                        # 取某个变量的 b、SE、P，不在模型里就给 NA
  out <- if (v %in% names(coef(fit))) coef(summary(fit))[v, c(1, 2, 4)] else rep(NA, 3)
  setNames(out, c("b", "SE", "P"))
}
round(t(sapply(mods, function(f) c(bmi = get(f, "bmi"), waist = get(f, "waist")))), 4)
\end{RCode}
\begin{ROut}
> round(t(sapply(mods, ...)), 4)
               bmi.b bmi.SE  bmi.P waist.b waist.SE waist.P
bmi           0.0634 0.0179 0.0005      NA       NA      NA
waist             NA     NA     NA  0.0266   0.0072  0.0003
age+bmi       0.0520 0.0138 0.0002      NA       NA      NA
age+waist         NA     NA     NA  0.0200   0.0056  0.0005
age+bmi+waist 0.0671 0.0552 0.2260 -0.0063   0.0223  0.7780
\end{ROut}
单独放都显著，一起放都不显著，waist 还变了号；但整个方程的 $R^2$ 和 $F$ 几乎不受影响。共线性伤的是“解释单个系数”，不是预测。

\subsection{VIF、容忍度与条件数}
\begin{RCode}[R · 共线性诊断]
vif(m)                                           # 1.03、16.02、16.14
1 / vif(m)                                       # 容忍度
1 / (1 - summary(lm(bmi ~ age + waist, data = d))$r.squared)   # 手算 bmi 的 VIF
sapply(x, function(v)                            # 一行算全部
  1 / (1 - summary(lm(reformulate(setdiff(x, v), v), data = d))$r.squared))
coef(summary(m))["bmi", 2] / coef(summary(lm(glu ~ age + bmi, d)))["bmi", 2]   # 4.00
sqrt(vif(m)["bmi"])                              # 4.00：√VIF 就是标准误膨胀倍数
ev <- eigen(cor(d[, x]))$values                  # 1.988、0.981、0.032
sqrt(max(ev) / min(ev))                          # 条件数 kappa = 7.93
kappa(scale(d[, x]), exact = TRUE)               # 同一个数
\end{RCode}
VIF 的经验界是 5（警惕）和 10（严重），条件数的界在教材里从 15 到 30 不等。本例 bmi 加进 waist 后标准误恰好放大 4.00 倍，等于 $\sqrt{16.02}$。

\begin{notebox}[课堂代码里的两个小问题]
\Verb|回归诊断.R| 末尾写的是 \Verb|vif(fit)|，而此时 \Verb|fit| 已被前一行改成了 \Verb|regsubsets| 对象，会报错 \Verb|缺少参数"id"|；\Verb|vif()| 只接受 lm 对象。另一行 \Verb|kappa(scale(d))| 把因变量 glu 也放进了矩阵，而且 \Verb|kappa()| 默认只给近似值。诊断共线性只放自变量，并写 \Verb|exact = TRUE|。
\end{notebox}

\subsection{岭回归：MASS::lm.ridge()}
\begin{RCode}[R · 拟合一串 lambda]
lam <- seq(0, 100, by = 0.5)
rr  <- lm.ridge(glu ~ age + bmi + waist, data = d, lambda = lam)
select(rr)                                  # HKB 1.16、L-W 1.21、GCV 最小在 7.5
rr$lambda[which.min(rr$GCV)]                # 7.5
round(t(rr$coef)[lam %in% c(0, 7.5, 20, 30, 60), ], 3)    # 画岭迹用的系数
round(coef(rr)[lam %in% c(0, 7.5, 30), ], 4)              # 写方程用的系数
\end{RCode}
\begin{ROut}
> round(t(rr$coef)[lam %in% c(0, 7.5, 20, 30, 60), ], 3)
        age   bmi  waist
  0.0 0.454 0.216 -0.050
  7.5 0.429 0.129  0.035
 20.0 0.397 0.100  0.061
 30.0 0.375 0.091  0.066
 60.0 0.322 0.079  0.070
> round(coef(rr)[lam %in% c(0, 7.5, 30), ], 4)
                age    bmi   waist
  0.0 2.4617 0.0314 0.0671 -0.0063
  7.5 2.3235 0.0297 0.0400  0.0044
 30.0 2.4546 0.0260 0.0283  0.0083
\end{ROut}

\begin{notebox}[lm.ridge 的两套系数]
\Verb|rr$coef| 是“变量 × λ”的矩阵，自变量已经标准化（按 $n$ 作分母的 SD），不含截距，画岭迹用它，画之前要 \Verb|t()| 转置；\Verb|coef(rr)| 是“λ × 变量”的矩阵，已还原到原始单位并带截距，写回归方程用它。$\lambda=0$ 那一行就是 OLS。
\end{notebox}

\begin{RCode}[R · 岭迹图]
matplot(rr$lambda, t(rr$coef), type = "l", lty = 1, lwd = 2,
        col = c("#0072B2", "#D55E00", "#009E73"),
        xlab = expression(lambda), ylab = "standardized coefficient")
abline(h = 0, col = "grey70"); abline(v = 7.5, lty = 2)
legend("topright", x, col = c("#0072B2", "#D55E00", "#009E73"), lty = 1, bty = "n")
plot(rr$lambda, rr$GCV, type = "l")         # GCV 曲线
\end{RCode}

\begin{figure}[htbp]
\centering
\includegraphics[width=0.95\textwidth]{figures/fig2_ridge.pdf}
\caption{datasetB 的岭迹与 GCV 曲线。waist 系数在 λ 很小时就穿过 0 变正，之后稳定；GCV 在 0–20 之间几乎是平的}
\label{fig:ridge}
\end{figure}

GCV 在 $\lambda=0$ 与 $\lambda=7.5$ 之间只差 0.6\%（图 \ref{fig:ridge}b）。GCV 衡量的是预测误差，而共线性不伤预测，所以它帮不上“读系数”的忙；岭参数要结合岭迹是否平稳、系数是否合理来定，并说明理由。

\subsection{自己写一遍岭回归}
\begin{RCode}[R · 按定义手算岭回归]
lambda <- 30
X  <- as.matrix(d[, x])
Xc <- scale(X, center = TRUE, scale = FALSE)
sc <- sqrt(colSums(Xc^2) / nrow(X))        # lm.ridge 用 n 作分母
Z  <- sweep(Xc, 2, sc, "/")                # 标准化后的自变量
b_std <- solve(crossprod(Z) + lambda * diag(3), crossprod(Z, d$glu - mean(d$glu)))
b_raw <- drop(b_std) / sc                  # 还原到原始单位
b0    <- mean(d$glu) - sum(b_raw * colMeans(X))
c(b0, b_raw)                               # 与 coef(lm.ridge(..., lambda = 30)) 完全相同
R <- cor(X); ry <- cor(X, d$glu)
solve(R + (lambda / nrow(X)) * diag(3), ry)   # 课件笔记的 (R + kI)^{-1} r_y，k = lambda/n
\end{RCode}
最后一行与课件笔记的相关阵形式一致，只是尺度不同：\Verb|lm.ridge| 的 $\lambda$ 对应 $k=\lambda/n$。不同软件对标准化的分母（$n$ 或 $n-1$）和 $\lambda$ 的归一方式不同，岭参数的数值不能跨软件直接比。

\subsection{比岭回归更常用的处理}
\begin{RCode}[R · 只留一个，或合成一个]
summary(lm(glu ~ age + bmi, data = d))                       # 只留 BMI
d$obesity <- (as.numeric(scale(d$bmi)) + as.numeric(scale(d$waist))) / 2
summary(lm(glu ~ age + obesity, data = d))                   # z 分数平均
pc <- prcomp(d[, c("bmi", "waist")], scale. = TRUE)
summary(pc)                                                  # 第一主成分占 98.4%
\end{RCode}
课件第 15 页给的顺序是：先由研究问题决定只留一个，其次合成一个指标，再次才是岭回归。合成指标与第一主成分在本例给出相同的 $t$（3.70），因为两个标准化变量的平均就是第一主成分的方向（主成分的正负号是任意的）。

\subsection{强影响点会改变岭参数的选择}
\begin{ROut}
> for (s in c("all", "del149")) ...
all : GCV 最优 lambda = 7.5 ；bmi、waist 同时为正的最小 lambda = 3
del149 : GCV 最优 lambda = 1.5 ；bmi、waist 同时为正的最小 lambda = 18.5
\end{ROut}
删掉 149 号后，要把 waist 拉成正号需要大得多的 λ。作业二第二题因此在“全样本、删 149、删 150、两人都删”四种样本下画岭迹，选了一个在四种情形下都已平稳的 λ = 30，而没有采用全样本 GCV 给出的 7.5（第 9 章）。

\section{一般线性模型：分类变量、方差分析与交互}
\srcnote{课件：第三讲第 2–19 页、第 25 页练习\quad 课堂代码：第三讲.R 前半、析因变异分解.docx\quad 脚本：\texttt{scripts/06\_一般线性模型\_分类变量与方差分析.R}\quad 数据：MASS::birthwt、glucose\_block、sbp\_factorial2x2、datasetC\_glm}

第三讲的核心一句话：t 检验、方差分析、回归、协方差分析都是 $y=X\beta+\varepsilon$，差别只在 $X$ 里放什么列。在 R 里这句话的意思是：它们都可以用 \Verb|lm()| 拟合，\Verb|t.test()|、\Verb|aov()| 只是换了一种输出格式。

\subsection{两样本 t 检验就是一个 0/1 列的回归}
\begin{RCode}[R · birthwt：母亲体重是否 > 110 磅]
library(MASS)
bw <- birthwt                                          # 189 名产妇，R 自带
bw$lwt_group <- cut(bw$lwt, breaks = c(-Inf, 110, Inf))
t.test(bwt ~ lwt_group, data = bw, var.equal = TRUE)    # 方差齐的成组 t 检验
fit_t <- lm(bwt ~ lwt_group, data = bw)
summary(fit_t)
c(t2 = coef(summary(fit_t))[2, "t value"]^2, F = anova(fit_t)[1, "F value"])
head(model.matrix(fit_t))                              # 一列 1，一列 0/1
\end{RCode}
\begin{ROut}
> summary(fit_t)
                    Estimate Std. Error t value Pr(>|t|)
(Intercept)          2656.34      97.31  27.298  < 2e-16 ***
lwt_group(110, Inf]   400.58     114.71   3.492 0.000598 ***
> c(t2 = ..., F = ...)
      t2        F
12.19377 12.19377
\end{ROut}
截距 2656.34 是参照组（≤ 110 磅）的平均出生体重，系数 400.58 是两组均数差，$t=3.492$、$P=0.0006$ 与 t 检验完全相同，$t^2=F$。\Verb|t.test()| 默认是 Welch 检验（\Verb|var.equal = FALSE|），要和回归对上必须写 \Verb|var.equal = TRUE|。

\subsection{分类变量不能当数字用}
\begin{RCode}[R · 四组：编成 1/2/3/4 与哑变量的区别]
bw$lwt4 <- cut(bw$lwt, breaks = c(-Inf, 110, 130, 150, Inf))
coef(lm(bwt ~ as.integer(lwt4), data = bw))   # 只有一个斜率 151.4：强加了等距、线性
coef(lm(bwt ~ lwt4, data = bw))               # 3 个哑变量，各组均数自由
\end{RCode}
\begin{ROut}
> coef(lm(bwt ~ lwt4, data = bw))
   (Intercept)  lwt4(110,130]  lwt4(130,150] lwt4(150, Inf]
     2656.3396       353.8318       428.3604       468.3270
\end{ROut}
因子进入 \Verb|lm()| 时，R 自动把 $k$ 个水平拆成 $k-1$ 个 0/1 列（处理编码，contr.treatment），系数 = 各组均数 − 参照组均数。这就是课件第 11 页的 $\hat y=2656.34+353.83g_2+428.36g_3+468.33g_4$。如果分组变量在数据里是 1、2、3 这样的数字（如 \Verb|glucose_block.csv| 的 subject），建模前一定要 \Verb|factor()|。

\subsection{单因素方差分析、参照组与多重比较}
\begin{RCode}[R · 整体检验、换参照组、两两比较]
fit4 <- lm(bwt ~ lwt4, data = bw)
anova(fit4)                                       # F = 4.26，P = 0.006
bw$lwt4_r <- relevel(bw$lwt4, ref = "(150, Inf]") # 名字要和 levels() 一字不差
coef(lm(bwt ~ lwt4_r, data = bw))                 # 系数全变
anova(lm(bwt ~ lwt4_r, data = bw))                # F、P 不变
TukeyHSD(aov(bwt ~ lwt4, data = bw))              # Tukey HSD
pairwise.t.test(bw$bwt, bw$lwt4, p.adjust.method = "bonferroni")
library(emmeans); pairs(emmeans(fit4, ~ lwt4), adjust = "tukey")
\end{RCode}
换参照组只改变“读哪一组对比”，不改变拟合值、残差、$R^2$ 和整体 $F$，这正是课件第 25 页练习 2 要说明的。多重比较三种写法中，\Verb|TukeyHSD()| 只接受 \Verb|aov()| 对象；\Verb|emmeans()| 可以接任何 \Verb|lm|，含协变量、含交互时也能用，后面几章一直用它。

\subsection{用 AIC 比较不嵌套的模型}
\begin{RCode}[R · 连续、二分、四分]
c(continuous  = AIC(lm(bwt ~ lwt, data = bw)),    # 3026.5
  two_groups  = AIC(fit_t),                        # 3021.2
  four_groups = AIC(fit4))                         # 3024.5
\end{RCode}
三个模型互不嵌套，不能用 \Verb|anova()| 比，但可以比 AIC：同一个因变量、同一批 189 人、都是 \Verb|lm()| 的最大似然，三条前提都满足。

\subsection{对比编码：处理编码与和为零编码}
\begin{RCode}[R · contrasts]
contrasts(bw$lwt4)                                    # 默认处理编码
options(contrasts = c("contr.sum", "contr.poly"))     # 改为和为零编码
coef(lm(bwt ~ lwt4, data = bw))       # 截距 = 各组均数的平均 2968.97；系数 = 效应 alphai
options(contrasts = c("contr.treatment", "contr.poly"))  # 用完立即恢复
\end{RCode}
\Verb|options(contrasts = ...)| 是全局设置，改了不恢复，后面所有模型的系数含义都会变。只想对一个模型改编码，可以写 \Verb|lm(y ~ g, contrasts = list(g = "contr.sum"))|。

\subsection{随机区组设计}
\begin{RCode}[R · 处理 + 区组]
gb <- read.csv("data/glucose_block.csv")      # 8 人 × 4 个放置时间
gb$time <- factor(gb$time_min); gb$subject <- factor(gb$subject)   # 必须是因子
fit_rb <- lm(glucose ~ time + subject, data = gb)
anova(fit_rb)
anova(lm(glucose ~ time, data = gb))          # 不控制区组
\end{RCode}
\begin{ROut}
> anova(fit_rb)
          Df Sum Sq Mean Sq F value    Pr(>F)
time       3 943.59 314.531  78.487 1.412e-11 ***
subject    7 806.22 115.174  28.740 2.199e-09 ***
Residuals 21  84.16   4.007
> anova(lm(glucose ~ time, data = gb))
time       3 943.59 314.531  9.8912 0.0001295 ***
Residuals 28 890.38  31.799
\end{ROut}
与课件第 15 页一致：$F=78.49$，$P=1.4\times10^{-11}$。不放区组时，区组间变异 806.22 全部进了误差，$F$ 从 78.5 掉到 9.9。\Verb|coef(fit_rb)| 给出课件第 14 页的 $b$（截距 96.03 是“第 1 人、0 分钟”的拟合值，time45 = −1.50 等是与 0 分钟之差），\Verb|model.matrix(fit_rb)| 就是那张 32 × 11 的设计矩阵。

\subsection{2×2 析因设计与交互作用}
\begin{RCode}[R · 交互模型、单元格均数、交互图]
f2 <- read.csv("data/sbp_factorial2x2.csv")
f2$a <- factor(f2$a); f2$b <- factor(f2$b)
fit_ab <- lm(sbp_drop ~ a * b, data = f2)
anova(fit_ab)
with(f2, tapply(sbp_drop, list(A = a, B = b), mean))      # 单元格均数
with(f2, interaction.plot(a, b, sbp_drop, type = "b"))    # 基础 R 交互图
\end{RCode}
\begin{ROut}
> anova(fit_ab)
          Df  Sum Sq Mean Sq F value    Pr(>F)
a          1 1278.24 1278.24 755.345 < 2.2e-16 ***
b          1  583.11  583.11 344.576 < 2.2e-16 ***
a:b        1   91.03   91.03  53.789 3.713e-09 ***
Residuals 44   74.46    1.69
> with(f2, tapply(sbp_drop, list(A = a, B = b), mean))
   B
A           0      1
  0  2.933333  7.150
  1 10.500000 20.225
\end{ROut}

\begin{figure}[htbp]
\centering
\includegraphics[width=0.5\textwidth]{figures/fig3_factorial.pdf}
\caption{2×2 析因的交互图（均数 ± 标准误）。两条线不平行：B 药存在时 A 药的降压效果更大}
\label{fig:factorial}
\end{figure}

交互 $F=53.79$，$P=3.7\times10^{-9}$，与课件第 19 页一致。课堂代码用 ggplot 画交互图，写法见脚本 §6.7；关键是 \Verb|stat_summary(fun = mean, geom = "line")| 加 \Verb|group = b|。

\subsection{手算析因的平方和分解}
课堂上的 \Verb|析因变异分解.docx| 列出了单元格均数、边际均数，以及“如果没有交互”时各格应有的均数。这几步在 R 里是：

\begin{RCode}[R · SSA、SSB、SSAB]
y <- f2$sbp_drop; gm <- mean(y); nc <- 12
cm <- with(f2, tapply(sbp_drop, list(a, b), mean))    # 2×2 单元格均数
am <- rowMeans(cm); bm <- colMeans(cm)                 # 边际均数
SSA  <- 2 * nc * sum((am - gm)^2)                      # 1278.235
SSB  <- 2 * nc * sum((bm - gm)^2)                      #  583.110
SSAB <- nc * sum((cm - outer(am, rep(1, 2)) - outer(rep(1, 2), bm) + gm)^2)   # 91.025
SSE  <- sum((y - gm)^2) - SSA - SSB - SSAB             #   74.459
outer(am, rep(1, 2)) + outer(rep(1, 2), bm) - gm       # 无交互时的格均数
\end{RCode}
最后一行给出 1.556、8.527、11.877、18.848，就是 docx 第三张表；实际格均数与它的差（如 $2.933-1.556$）就是交互效应 $(\alpha\beta)_{ij}$。\Verb|outer(u, v)| 生成 $u_iv_j$ 矩阵，\Verb|outer(am, rep(1, 2))| 就是把行边际均数铺满每一列。作业三的 3×3 手算用的是同一套写法。

\subsection{III 型平方和必须配和为零编码}
\label{sec:type3}
\begin{RCode}[R · Anova(type = "III") 的正确与错误写法]
library(car)
options(contrasts = c("contr.sum", "contr.poly"))
Anova(lm(sbp_drop ~ a * b, data = f2), type = "III")   # 正确
options(contrasts = c("contr.treatment", "contr.poly"))
Anova(lm(sbp_drop ~ a * b, data = f2), type = "III")   # 错误：主效应那两行不对
\end{RCode}
\begin{ROut}
> Anova(..., type = "III")      # contr.sum
a           1278.2  1  755.345 < 2.2e-16 ***
b            583.1  1  344.576 < 2.2e-16 ***
a:b           91.0  1   53.789 3.713e-09 ***
> Anova(..., type = "III")      # 默认 contr.treatment
a           343.53  1 203.000 < 2.2e-16 ***
b           106.68  1  63.041 4.922e-10 ***
a:b          91.03  1  53.789 3.713e-09 ***
\end{ROut}
在默认处理编码下，含交互模型里 a 的系数是“b 取参照水平时”A 的效应，III 型检验的就是它，所以 SS 变成 343.5。第三讲课堂代码在 \Verb|Anova(fit, type = "III")| 之前写了 \Verb|options(contrasts = c("contr.sum", "contr.poly"))|，就是这个原因。平衡设计里 I、II、III 型相同；不平衡时三者不同，下一节的练习 3 就是例子。

\subsection{交互显著之后：简单效应}
\begin{RCode}[R · emmeans 的条件比较]
pairs(emmeans(fit_ab, ~ a | b))     # 固定 B 的每个水平，比较 A
pairs(emmeans(fit_ab, ~ b | a))     # 固定 A 的每个水平，比较 B
contrast(emmeans(fit_ab, ~ a * b), interaction = "pairwise")   # 交互对比本身
\end{RCode}
\begin{ROut}
> pairs(emmeans(fit_ab, ~ a | b))
b = 0:
 contrast estimate    SE df t.ratio p.value
 a0 - a1     -7.57 0.531 44 -14.248 <0.0001
b = 1:
 a0 - a1    -13.07 0.531 44 -24.620 <0.0001
\end{ROut}
不用 B 药时 A 药多降 7.57 mmHg，用 B 药时多降 13.07 mmHg；两者之差 5.51 就是交互对比的估计值。\Verb!~ a | b! 里竖线读作“在 b 的每个水平下”，第 8 章重复测量的两个方向的简单效应也是这样写。

\subsection{课件第 25 页四道练习}
数据 \Verb|datasetC_glm.csv|，完整代码在脚本 §6.11，这里只列要点和结果。

\begin{exercisebox}{glu \textasciitilde{} sex：t 检验与回归逐项核对}{glm1}
\Verb|t.test(glu ~ sex, data = g, var.equal = TRUE)| 与 \Verb|lm(glu ~ sex, data = g)|：均数差 0.4895，$t=4.641$，$P=5.7\times10^{-6}$，$t^2=F=21.54$。设计矩阵前 3 行 \Verb|model.matrix(m1)[1:3, ]| 都是 (1, 1)，因为前三人都是男性、female 是参照组。
\end{exercisebox}

\begin{exercisebox}{glu \textasciitilde{} bmi3：换参照组}{glm2}
以 normal 为参照：截距 4.838，overwt +0.843，obese +1.203；以 obese 为参照：截距 6.041，normal −1.203，overwt −0.360。整体 $F=50.72$ 不变，\Verb|all.equal(fitted(m2a), fitted(m2b))| 为 TRUE。
\end{exercisebox}

\begin{exercisebox}{glu \textasciitilde{} sex * ov：不平衡的两因素设计}{glm3}
性别与超重的人数不平衡（61/59/39/81），三种平方和不同：sex 一行 I 型 14.38、II 型 6.34、III 型 4.15；交互一行三者都相同，$F=12.99$，$P=0.0004$。简单效应：女性中超重者血糖高 0.557，男性中高 1.216（\Verb!pairs(emmeans(m3, ~ ov | sex))!），男性中超重的影响更大。
\end{exercisebox}

\begin{exercisebox}{glu \textasciitilde{} bmi3 + age：协方差分析}{glm4}
见下一章第 \ref{sec:ancova-c} 节。
\end{exercisebox}

\section{协方差分析}
\srcnote{课件：第三讲第 20–24 页；第四讲第 3–13 页\quad 课堂代码：第三讲.R 后半、第四讲.R 前半\quad 脚本：\texttt{scripts/07\_协方差分析.R}\quad 数据：datasetC\_glm、bp、knee}

协方差分析就是“分组因子 + 一个连续协变量”的线性模型。R 里的流程固定为三步：

\begin{definitionbox}{协方差分析的三步}{ancova}
\begin{enumerate}
\item \textbf{斜率同质性}：\Verb|anova(lm(y ~ g + x), lm(y ~ g * x))|，P > 0.05 才继续。
\item \textbf{校正模型}：\Verb|fit <- lm(y ~ x + g)|，用 \Verb|Anova(fit, type = "II")| 或把协变量写在前面的 \Verb|anova(fit)| 检验组别。
\item \textbf{校正均数与两两比较}：\Verb|emm <- emmeans(fit, ~ g)|，\Verb|pairs(emm, adjust = "bonferroni")|。
\end{enumerate}
\end{definitionbox}

\subsection{第三讲的例子：三组 BMI，年龄不齐}
\label{sec:ancova-c}
\begin{RCode}[R · 为什么要校正]
g <- read.csv("data/datasetC_glm.csv")
g$bmi3 <- factor(g$bmi3, levels = c("normal", "overwt", "obese"))   # 按专业顺序排水平
aggregate(cbind(age, glu) ~ bmi3, data = g, FUN = mean)
anova(lm(age ~ bmi3, data = g))            # 三组年龄不同：F = 8.76，P = 0.0002
\end{RCode}
\begin{ROut}
> aggregate(cbind(age, glu) ~ bmi3, data = g, FUN = mean)
    bmi3      age      glu
1 normal 49.30303 4.838081
2 overwt 54.56250 5.681339
3  obese 56.44828 6.041379
\end{ROut}
\Verb|read.csv()| 读进来的字符变量按字母顺序排水平（normal、obese、overwt），课堂代码就是这个顺序；这里用 \Verb|factor(..., levels = )| 改成专业上的顺序，参照组仍是 normal，结果相同，只是输出更好读。

\paragraph{第一步：斜率同质性。}
\begin{RCode}[R · 平行性检验]
sapply(split(g, g$bmi3), function(s) coef(lm(glu ~ age, data = s))["age"])
m_par <- lm(glu ~ bmi3 + age, data = g)      # 共同斜率
m_int <- lm(glu ~ bmi3 * age, data = g)      # 各组斜率不同
anova(m_par, m_int)
\end{RCode}
\begin{ROut}
> sapply(...)
normal.age overwt.age  obese.age
0.03503285 0.03244822 0.02993767
> anova(m_par, m_int)
  Res.Df    RSS Df Sum of Sq      F Pr(>F)
1    236 92.498
2    234 92.419  2  0.078634 0.0995 0.9053
\end{ROut}
三组斜率 0.035、0.032、0.030，$F=0.0995$，$P=0.905$，可以用共同斜率。与课件第 21 页一致。

\paragraph{第二步：校正模型——协变量要写在前面。}
\label{sec:ancova-order}
\begin{RCode}[R · 三种写法对比]
anova(lm(glu ~ age + bmi3, data = g))     # 协变量在前：bmi3 一行是校正年龄后的
anova(lm(glu ~ bmi3 + age, data = g))     # 分组在前：bmi3 一行没有校正年龄
library(car); Anova(m_par, type = "II")   # 与顺序无关，推荐
\end{RCode}
\begin{ROut}
> anova(lm(glu ~ age + bmi3, data = g))
           Df Sum Sq Mean Sq F value    Pr(>F)
age         1 49.976  49.976 127.508 < 2.2e-16 ***
bmi3        2 30.769  15.385  39.252 1.922e-15 ***
> anova(lm(glu ~ bmi3 + age, data = g))
bmi3        2 51.927 25.9636  66.244 < 2.2e-16 ***
age         1 28.817 28.8173  73.525 1.327e-15 ***
> Anova(m_par, type = "II")
          Sum Sq  Df F value    Pr(>F)
bmi3      30.769   2  39.252 1.922e-15 ***
age       28.817   1  73.525 1.327e-15 ***
\end{ROut}
这是第 2 章“I 型平方和看顺序”在协方差分析里的后果：分组写在前面时，bmi3 那一行（$F=66.24$）就是不校正年龄的单因素方差分析。协方差分析要检验的是校正后的组间差异（$F=39.25$），所以要么把协变量写在前面，要么用 \Verb|Anova(type = "II")|。课堂代码 \Verb|anova(lm(glu ~ bmi3 + age + bmi3*age))| 和第四讲的 \Verb|anova(lm(post ~ group + base))| 都是分组在前，见第 \ref{sec:bp} 节。

\paragraph{校正前后对照。}
\begin{RCode}[R · 粗差与校正差]
m_crude <- lm(glu ~ bmi3, data = g)
rbind(crude = coef(summary(m_crude))[2:3, 1:2], adjusted = coef(summary(m_par))[2:3, 1:2])
c(sigma(m_crude), sigma(m_par)); c(summary(m_crude)$r.squared, summary(m_par)$r.squared)
\end{RCode}
\begin{table}[htbp]
\centering\small
\caption{datasetC\_glm：校正年龄前后与正常组的差值}
\begin{tabular}{@{}lS[table-format=1.3]S[table-format=1.3]S[table-format=1.3]S[table-format=1.3]@{}}
\toprule
 & \multicolumn{2}{c}{超重 vs 正常} & \multicolumn{2}{c}{肥胖 vs 正常}\\
\cmidrule(lr){2-3}\cmidrule(lr){4-5}
 & {差值} & {标准误} & {差值} & {标准误}\\
\midrule
未校正 \Verb|lm(glu ~ bmi3)| & 0.843 & 0.099 & 1.203 & 0.151\\
校正年龄 \Verb|lm(glu ~ bmi3 + age)| & 0.669 & 0.089 & 0.966 & 0.135\\
\bottomrule
\end{tabular}
\TableNote{剩余标准差 0.715 → 0.626，$R^2$ 0.300 → 0.466。}
\end{table}
差值缩小两成（控混杂），标准误同时变小（提精度），与课件第 22 页一致。

\paragraph{第三步：校正均数。}
\begin{RCode}[R · 三种算法]
library(emmeans)
emm <- emmeans(m_par, ~ bmi3); emm                         # (1) emmeans
predict(m_par, newdata = data.frame(bmi3 = levels(g$bmi3), age = mean(g$age)))   # (2) 在总平均年龄处预测
lxy <- sum(sapply(split(g, g$bmi3), function(s)            # (3) 手算：合并组内斜率
  sum((s$age - mean(s$age)) * (s$glu - mean(s$glu)))))
lxx <- sum(sapply(split(g, g$bmi3), function(s) sum((s$age - mean(s$age))^2)))
b_w <- lxy / lxx                                           # 0.03316 = coef(m_par)["age"]
tapply(g$glu, g$bmi3, mean) - b_w * (tapply(g$age, g$bmi3, mean) - mean(g$age))
\end{RCode}
\begin{ROut}
> emm
 bmi3   emmean     SE  df lower.CL upper.CL
 normal   4.95 0.0642 236     4.82     5.07
 overwt   5.62 0.0596 236     5.50     5.73
 obese    5.91 0.1170 236     5.68     6.15
> tapply(...) - b_w * (...)
  normal   overwt    obese
4.948095 5.616956 5.914467
\end{ROut}
三种算法结果相同：校正均数就是“把三组都放到总平均年龄 52.6 岁时”模型的预测值，等于组均数减去 $b_w(\bar x_i-\bar x)$。手算那几行对应课堂代码里的 \Verb|分子 / 分母|（\Verb|cov(...)*98| 就是组内离均差积和，98 = 组内人数 − 1）。

\begin{RCode}[R · 两两比较]
pairs(emm, adjust = "tukey")                       # 课件第 24 页列的几种方法之一
pairs(emm, adjust = "bonferroni")
confint(pairs(emm, adjust = "bonferroni"))
contrast(emm, method = "trt.vs.ctrl", ref = "normal")   # 都与正常组比（Dunnett）
\end{RCode}
\begin{ROut}
> pairs(emm, adjust = "tukey")
 contrast        estimate     SE  df t.ratio p.value
 normal - overwt   -0.669 0.0887 236  -7.539 <0.0001
 normal - obese    -0.966 0.1350 236  -7.156 <0.0001
 overwt - obese    -0.298 0.1310 236  -2.277  0.0610
\end{ROut}
\Verb|adjust =| 可以取 \Verb|"tukey"|、\Verb|"bonferroni"|、\Verb|"sidak"|、\Verb|"scheffe"|、\Verb|"none"|；与对照组比较用 \Verb|method = "trt.vs.ctrl"|（Dunnett）。

\begin{figure}[htbp]
\centering
\includegraphics[width=0.98\textwidth]{figures/fig4_ancova.pdf}
\caption{(a) datasetC\_glm：三条平行线，竖虚线为总平均年龄，线在虚线处的高度就是校正均数；(b) bp.csv：校正基线后两组的高低顺序反转}
\label{fig:ancova}
\end{figure}

画平行线的思路：造一张“三组 × 一串年龄”的网格，用 \Verb|predict(m_par, newdata = grid)| 得到拟合值，再按组画线。脚本 §7.6 用 \Verb|expand.grid()| 加 ggplot 完成，这也是课堂代码的写法。

\paragraph{练习 4 后半：把年龄切成三组再校正。}\Verb|lm(glu ~ bmi3 + age3)| 得到的超重、肥胖差值是 0.681、1.030，比连续年龄的 0.669、0.966 大，剩余标准差也更大（0.641 vs 0.626）：切组后组内仍有年龄差，校正不彻底。协变量是连续的就保持连续。

\subsection{第四讲的例子：基线不齐，结论反转}
\label{sec:bp}
\begin{RCode}[R · bp.csv]
bp <- read.csv("data/bp.csv", fileEncoding = "UTF-8")
bp$group <- factor(bp$group, levels = c("方案A", "方案B"), labels = c("A", "B"))
aggregate(cbind(base, post) ~ group, data = bp, FUN = mean)   # A 158→137，B 142→131
t.test(post ~ group, data = bp, var.equal = TRUE)             # 只比终点：A 更高，P = 0.035
anova(lm(post ~ group * base, data = bp))                     # 交互 P = 0.929，平行
fit <- lm(post ~ base + group, data = bp)                     # 协变量在前
anova(fit)
anova(lm(post ~ group + base, data = bp))                     # 分组在前（课件的写法）
pairs(emmeans(fit, ~ group))
\end{RCode}
\begin{ROut}
> anova(fit)
          Df  Sum Sq Mean Sq F value    Pr(>F)
base       1 11627.8 11627.8 195.612 < 2.2e-16 ***
group      1  1506.8  1506.8  25.349 2.206e-06 ***
Residuals 97  5766.0    59.4
> anova(lm(post ~ group + base, data = bp))
group      1    841   841.0  14.148 0.0002891 ***
base       1  12294 12293.6 206.813 < 2.2e-16 ***
> pairs(emmeans(fit, ~ group))
 contrast estimate   SE df t.ratio p.value
 A - B       -9.42 1.87 97  -5.035 <0.0001
\end{ROut}
校正均数 A 129.2、B 138.6，A 方案低 9.42 mmHg（$t=-5.04$），与只比终点的结论相反。

\begin{notebox}[课件第 10–11 页方差分析表里的“分组”一行]
那两张表中分组 SS = 841.0、$F=14.148$、$P=0.00029$，来自 \Verb|anova(lm(post ~ group + base))|——分组写在基线前面，841.0 恰好是两组终点均数之差 5.8 对应的平方和（$50\times50/100\times5.8^2$），是\textbf{未校正}的。校正基线后的组间检验是 SS = 1506.8、$F(1,97)=25.35$、$P=2.2\times10^{-6}$，与 emmeans 两两比较的 $t^2=5.035^2$ 一致。课件给出的校正均数 129.17、138.59 和差值 −9.42 本身没有问题。写作业时请用协变量在前的 \Verb|anova()| 或 \Verb|Anova(type = "II")|。
\end{notebox}

\paragraph{前提检查。}\Verb|t.test(base ~ group)|：两组基线相差 15.8 mmHg（$t=6.79$），在随机化试验里这么大的基线差本身就值得核查；\Verb|anova(fit, update(fit, . ~ . + I(base^2)))| 检查线性（$P=0.40$）；\Verb|shapiro.test(resid(fit))|（$P=0.63$）；\Verb|leveneTest(resid(fit) ~ bp$group)|（$P=0.16$）；\Verb|outlierTest(fit)| 没有离群点。协变量必须是处理前测得、不受处理影响的变量，这一条靠设计保证，检验不出来。

\paragraph{比降幅行不行。}\Verb|t.test(I(base - post) ~ group)| 比较降幅，A 降 21.0、B 降 11.0。差值分析相当于假定基线对终点的斜率恰好为 1；本例协方差分析估计的斜率是 0.964，两种方法结论一致。斜率明显偏离 1 时，差值分析会把基线差异的影响混进来，协方差分析更稳妥。

\subsection{knee.csv：另一个反转的例子}
第四讲文件夹里的 \Verb|knee.csv| 是两组膝关节术后康复数据（课件第 13 页三组练习的两组版本）：未校正时 B 组术后评分高（77.2 vs 75.6）；平行性检验 $P=0.27$；校正术前评分后 A 组校正均数 78.4、B 组 74.5，A − B = 3.86，$P=0.036$。代码在脚本 §7.10。三组的版本 \Verb|knee3.csv| 是作业四第 1 题（第 9 章）。

\section{重复测量方差分析}
\srcnote{课件：第四讲第 14–35 页\quad 课堂代码：第四讲.R 后半\quad 脚本：\texttt{scripts/08\_重复测量方差分析.R}\quad 数据：dm（宽型）、df\_long（长型），40 例 2 型糖尿病 × 3 个时点 HbA1c}

\subsection{数据准备：长型，时间和个体都是因子}
\begin{RCode}[R · 宽转长]
library(tidyr); library(afex); library(emmeans)
dm <- read.csv("data/dm.csv", fileEncoding = "UTF-8")
dl <- pivot_longer(dm, cols = c(hba1c_0, hba1c_3, hba1c_6),
                   names_to = "time", values_to = "hba1c")
dl$time  <- factor(dl$time, levels = c("hba1c_0", "hba1c_3", "hba1c_6"),
                   labels = c("T0", "M3", "M6"))
dl$group <- factor(dl$group); dl$id <- factor(dl$id)
aggregate(hba1c ~ group + time, data = dl,
          FUN = function(x) round(c(mean = mean(x), sd = sd(x)), 2))
\end{RCode}
\begin{ROut}
> aggregate(hba1c ~ group + time, ...)
  group time hba1c.mean hba1c.sd
1 A药组   T0       8.21     0.68
2 B药组   T0       8.80     0.65
3 A药组   M3       7.01     0.62
4 B药组   M3       8.01     0.63
5 A药组   M6       6.66     0.81
6 B药组   M6       7.66     0.72
\end{ROut}
与课件第 24 页的表一致。第 31 页折线图上标的 B 组 8.62、7.62 和 A 组 7.62、6.62 与数据不符，以表格为准。

\begin{figure}[htbp]
\centering
\includegraphics[width=0.98\textwidth]{figures/fig5_rm.pdf}
\caption{dm.csv 两组 HbA1c 的轨迹。(a) 各时点均数与 95\% CI；(b) 每个人的轨迹（细线）与组均数（粗线）}
\label{fig:rm}
\end{figure}

\subsection{单组：只有时间一个因素}
\begin{RCode}[R · aov\_ez 单组]
dA <- droplevels(subset(dl, group == "A药组"))     # 20 人 × 3 时点
mod1 <- aov_ez(id = "id", dv = "hba1c", data = dA, within = "time")
summary(mod1)                                      # 一次给出 F、Mauchly、GG、HF
anova(mod1, correction = "GG")                     # 按 GG 校正后的表
summary(aov(hba1c ~ time + Error(id / time), data = dA))   # 基础 R 的等价写法
\end{RCode}
\begin{ROut}
> summary(mod1)
            Sum Sq num Df Error SS den Df F value    Pr(>F)
(Intercept) 3191.6      1   25.344     19  2392.7 < 2.2e-16 ***
time          26.7      2    3.249     38   156.4 < 2.2e-16 ***

Mauchly Tests for Sphericity
     Test statistic p-value
time        0.80165 0.13674

Greenhouse-Geisser and Huynh-Feldt Corrections
      GG eps Pr(>F[GG])
time 0.83448  2.821e-16 ***
       HF eps   Pr(>F[HF])
time 0.905294 1.802841e-17
\end{ROut}
读法：time 一行 $F(2,38)=156.4$；Error SS 3.249 是受试者内误差，受试者间变异 25.344（课件第 18 页的“受试者”一行）显示在 (Intercept) 行的 Error SS 里，不要把 Intercept 那一行的 Sum Sq 3191.6 当成变异来源；Mauchly $W=0.802$，$P=0.137$；GG $\varepsilon=0.834$、HF $\varepsilon=0.905$。与课件第 18–19 页一致。

\begin{notebox}[aov\_ez 的参数]
\Verb|id| 是个体编号列，\Verb|dv| 是因变量列，\Verb|within| 是组内因素（时间），\Verb|between| 是组间因素（分组），都写成字符串。数据必须是长型，每人每个时点一行。afex 会提示 “Contrasts set to contr.sum”，它内部用和为零编码算 III 型平方和，不影响你在别处的设置。课件笔记用的 \Verb|ez::ezANOVA(data, dv = y, wid = id, within = t)| 是另一个包的同类函数，参数不加引号，结果相同（附录 B）。
\end{notebox}

单组的变异分解可以手算核对：

\begin{RCode}[R · 变异分解手算（单组）]
Y  <- as.matrix(dm[dm$group == "A药组", c("hba1c_0", "hba1c_3", "hba1c_6")])
n  <- nrow(Y); k <- ncol(Y); gm <- mean(Y)
SS_time    <- n * sum((colMeans(Y) - gm)^2)          # 26.744
SS_subject <- k * sum((rowMeans(Y) - gm)^2)          # 25.344
SS_error   <- sum((Y - gm)^2) - SS_time - SS_subject # 3.249
(SS_time / (k - 1)) / (SS_error / ((n - 1) * (k - 1)))   # F = 156.4
\end{RCode}
受试者间变异被单独剥离出去，误差只剩同一个人内部的波动，所以 $F$ 很大。这与随机区组（第 6.6 节）的分解完全相同，区组就是“人”。

\subsection{球形检验与 GG、HF 校正：手算一遍}
\begin{RCode}[R · 差值方差、Mauchly W、eps]
S <- cov(Y); round(S, 3)                            # 课件第 21 页的协方差阵
D <- cbind(Y[, 1] - Y[, 2], Y[, 1] - Y[, 3], Y[, 2] - Y[, 3])
round(apply(D, 2, var), 3)                          # 0.117 0.245 0.151：球形要求三者相等
C  <- rbind(c(1, -1, 0) / sqrt(2), c(1, 1, -2) / sqrt(6))   # 正交单位化对比，k-1 行
Sc <- C %*% S %*% t(C); p <- nrow(Sc); nu <- n - 1
W   <- det(Sc) / (sum(diag(Sc)) / p)^p              # 0.8017
chi <- -(nu - (2 * p^2 + p + 2) / (6 * p)) * log(W) # 3.979
pchisq(chi, p * (p + 1) / 2 - 1, lower.tail = FALSE)   # 0.1367
lam <- eigen(Sc)$values
sum(lam)^2 / (p * sum(lam^2))                       # GG：0.8345
mauchly.test(lm(Y ~ 1), X = ~ 1)                    # 基础 R：W = 0.80165，P = 0.1367
\end{RCode}
几点说明：

\begin{itemize}
\item $S_c$ 必须用\textbf{正交、单位化}的 $k-1$ 个对比去变换，课堂代码那段“手动逐行标准化”的循环做的就是单位化；直接用三个两两差值的协方差阵，秩只有 2，行列式为 0。
\item 卡方近似的校正项按 $(2p^2+p+2)/(6p)$ 计算（$p=k-1$），这与 \Verb|mauchly.test()|、afex 一致；课件第 22 页写的 $(2p^2+2p+3)/(6p)$ 与软件对不上。
\item HF 系数的单组公式 \Verb|(n*p*eps - 2) / (p*(nu - p*eps))|（上限 1）在脚本 §8.5 里，结果 0.905。
\item 校正就是把分子、分母自由度都乘以 $\varepsilon$，$F$ 值不变：\Verb|pf(156.4, 2*0.834, 38*0.834, lower.tail = FALSE)| 得到 GG 校正的 P。
\item $W$ 的 $P=0.137$ 只说明没有足够证据拒绝球形，差值方差实际相差两倍；$\varepsilon$ 只有 0.834，报告时可以同时给出 GG 校正结果（课件第 19、21 页）。
\end{itemize}

\subsection{多组：药物（组间）× 时间（组内）}
\begin{RCode}[R · aov\_ez 多组]
mod <- aov_ez(id = "id", dv = "hba1c", data = dl, within = "time", between = "group")
summary(mod)
anova(mod, correction = "GG", es = "pes")                       # 带偏 eta^2
summary(aov(hba1c ~ group * time + Error(id / time), data = dl))  # 两层误差项
\end{RCode}
\begin{ROut}
> summary(mod)
            Sum Sq num Df Error SS den Df  F value    Pr(>F)
(Intercept) 7159.5      1   46.313     38 5874.486 < 2.2e-16 ***
group         22.3      1   46.313     38   18.276 0.0001238 ***
time          39.5      2    7.529     76  199.550 < 2.2e-16 ***
group:time     1.1      2    7.529     76    5.588 0.0054428 **

Mauchly Tests for Sphericity
           Test statistic p-value
time              0.89178 0.12017
group:time        0.89178 0.12017
            GG eps Pr(>F[GG])
time       0.90235    < 2e-16 ***
group:time 0.90235    0.00722 **
\end{ROut}
三个 $F$ 与课件第 27 页一致：药物 $F(1,38)=18.28$，时间 $F(2,76)=199.55$，交互 $F(2,76)=5.59$，$P=0.005$（GG 校正后 0.007）。注意分母不同：药物用受试者间误差（Error SS 46.313，df 38），时间和交互用受试者内误差（7.529，df 76）。\Verb|aov(... + Error(id/time))| 的输出把这两层分成 \Verb|Error: id| 与 \Verb|Error: id:time| 两块，结构一目了然。

\begin{notebox}[错误示范：当作独立观测]
\Verb|anova(lm(hba1c ~ group * time, data = dl))| 把 120 个观测当作 120 个人：误差自由度变成 114，组间 $F$ 被夸大到 47.2，而交互 $F$ 只有 1.17（$P=0.31$）——个体差异全混进了误差，真实存在的交互反而检验不出来。同理，在每个时点分别做 t 检验也是错的。
\end{notebox}

\paragraph{多组的变异分解手算。}脚本 §8.7 按课件第 29–30 页的五个公式逐项计算：组间 22.2741、受试者间 46.3125、时间 39.5372、交互 1.1072、受试者内误差 7.5290，合计 116.76，三个 $F$ 为 18.28、199.55、5.59。代码里 \Verb|mik[grp, ]| 把“组 × 时点”均数按每个人所在的组展开成 40 × 3 矩阵，\Verb|outer(mij, rep(1, K))| 把个体均数铺成 40 × 3，然后按公式加减即可。

\begin{notebox}[易错：tapply 返回的是一维数组]
\Verb|mi <- tapply(rowMeans(Y), grp, mean)| 返回的是带 dim 属性的一维数组，\Verb|mi[grp]| 与 40 × 3 的矩阵相加会报“非整合陈列”（non-conformable arrays）。写成 \Verb|as.vector(tapply(...))| 就好了。
\end{notebox}

\paragraph{多组时的球形检验}用的是合并组内协方差阵 $S_p=\sum_i(n_i-1)S_i/(N-a)$，自由度 $\nu=N-a$；其余步骤与单组相同，脚本 §8.8 算出 $W=0.892$、$P=0.120$、GG $\varepsilon=0.902$，与 afex 一致。课件笔记例 12-3（三组麻醉诱导，附录 B）也是这个算法。

\subsection{交互显著之后：两个方向的简单效应}
\begin{RCode}[R · emmeans]
pairs(emmeans(mod, ~ time | group), adjust = "bonferroni")   # (1) 组内：每组各时点比
pairs(emmeans(mod, ~ group | time))                          # (2) 组间：每个时点两组比
contrast(emmeans(mod, ~ time * group), interaction = c("consec", "pairwise"))
\end{RCode}
\begin{ROut}
> pairs(emmeans(mod, ~ time | group), adjust = "bonferroni")
group = A药组:
 contrast estimate     SE df t.ratio p.value
 T0 - M3     1.205 0.0818 38  14.725 <0.0001
 T0 - M6     1.560 0.1050 38  14.858 <0.0001
 M3 - M6     0.355 0.1100 38   3.241  0.0074
group = B药组:
 T0 - M3     0.800 0.0818 38   9.776 <0.0001
 T0 - M6     1.150 0.1050 38  10.953 <0.0001
 M3 - M6     0.350 0.1100 38   3.195  0.0084
> pairs(emmeans(mod, ~ group | time))
time = T0:   A药组 - B药组   -0.590 0.210 38  -2.808  0.0078
time = M3:   A药组 - B药组   -0.995 0.197 38  -5.058 <0.0001
time = M6:   A药组 - B药组   -1.000 0.243 38  -4.123  0.0002
> contrast(..., interaction = c("consec", "pairwise"))
 time_consec group_pairwise estimate    SE df t.ratio p.value
 M3 - T0     A药组 - B药组    -0.405 0.116 38  -3.500  0.0012
 M6 - M3     A药组 - B药组    -0.005 0.155 38  -0.032  0.9744
\end{ROut}
与课件第 32 页一致：两组各阶段都显著下降；两组差距在 0–3 个月拉开 0.405，3–6 个月不再变化。最后一个交互对比把“差距在哪个阶段变化”直接检验出来，是解释交互作用最清楚的一步。两个方向的多重比较要分开校正：组内每组 $k(k-1)/2$ 次，组间每个时点 $g(g-1)/2$ 次。

\paragraph{不依赖球形假设的多元检验。}\Verb|summary(mod$Anova, multivariate = TRUE)| 给出 Pillai、Wilks 等多元检验，交互 $F(2,37)=7.36$，$P=0.002$，与单变量检验的结论一致。

\begin{resultbox}[重复测量的报告要点]
写明球形检验结果和采用的校正方法（未校正、GG 或 HF）；三个 $F$ 都给出统计量、自由度、P 值和效应量（\Verb|es = "pes"| 或 \Verb|"ges"|）；交互显著时以简单效应描述“哪一组、在哪个阶段”变化不同；配一张均数轨迹图。
\end{resultbox}

\section{四次作业完整题解与 R 代码}
\srcnote{脚本：\texttt{homework/hw1–hw4\_*.R}，均已在学习包内重新跑通。本章按“题目要求 → 关键 R 语句 → 统计结果与图表 → 结果解读与局限”的统一体例完整收录四次作业的全部题解，方便复习时把理论方法与实战代码一一对照。}

% ============================================================================
\subsection{作业一：27 例糖尿病患者空腹血糖的多重线性回归}

数据为 \texttt{diabetes27\_hw.csv}（$n=27$，无缺失值），包含总胆固醇 \texttt{tc}（mmol/L）、甘油三酯 \texttt{tg}（mmol/L）、空腹胰岛素 \texttt{insulin}（mU/L）、糖化血红蛋白 \texttt{hba1c}（\%）和空腹血糖 \texttt{glucose}（mmol/L）。读入数据后先核对描述统计与简单相关：

\begin{RCode}[作业一：读入与描述统计]
d <- read.csv("data/diabetes27_hw.csv")
round(sapply(d[, -1], mean), 2); round(sapply(d[, -1], sd), 2)
round(cor(d[, -1]), 3)
\end{RCode}

\begin{table}[htbp]
\centering\small
\caption{作业一各变量的描述统计及其与空腹血糖的简单相关（$n=27$）}
\begin{tabular}{@{}llccc@{}}
\toprule
变量 & 含义（单位） & 均数 & 标准差 & 与 glucose 的 $r$\\
\midrule
tc      & 总胆固醇（mmol/L） & 5.81  & 1.59 & 0.559\\
tg      & 甘油三酯（mmol/L） & 2.84  & 2.57 & 0.459\\
insulin & 空腹胰岛素（mU/L） & 6.15  & 3.67 & $-0.510$\\
hba1c   & 糖化血红蛋白（\%） & 9.12  & 1.82 & 0.610\\
glucose & 空腹血糖（mmol/L） & 11.93 & 2.93 & ---\\
\bottomrule
\end{tabular}
\TableNote{自变量间相关最强的是 tc 与 tg（$r=0.632$），其次是 tc 与 hba1c（$0.415$）、tc 与 insulin（$-0.355$）。25 号患者的 tc（11.54）、tg（10.89）与 glucose（20.0）均为全样本最高。}
\end{table}

\subsubsection*{第 1 题\quad 初始模型与 drop1 偏 $F$ 检验}
\emph{题目要求：根据专业知识建立含 4 个自变量的初始模型后，用 \texttt{drop1} 做偏 $F$ 检验确定最终模型，写出检验过程（$F$ 值、自由度与 $P$ 值）。}

\paragraph{1. 初始模型与检验原理。}糖化血红蛋白反映近 2--3 个月平均血糖，预期正相关；空腹胰岛素反映胰岛 $\beta$ 细胞分泌功能，预期负相关；甘油三酯与总胆固醇同属脂代谢指标，胰岛素抵抗时血脂与血糖常伴随升高。因此四变量均纳入初始模型：
\[
\text{glucose}=\beta_0+\beta_1\,\text{tc}+\beta_2\,\text{tg}+\beta_3\,\text{insulin}+\beta_4\,\text{hba1c}+\varepsilon .
\]
全模型 $\mathrm{SSE}=88.841$（$\mathrm{df}=22$），$R^2=0.601$，调整 $R^2=0.528$，整体 $F(4,22)=8.28$，$P<0.001$。对模型中单个变量 $x_j$ 检验 $H_0:\beta_j=0$ 的偏 $F$ 统计量为 $F_j=(\mathrm{SSE}_{(-j)}-\mathrm{SSE})/\mathrm{MSE}\sim F(1,n-p-1)$，等价于 $t_j^2$。后退法每步只删 $F$ 最小且 $P\ge 0.05$ 的一个变量，直至全部留下变量的 $P<0.05$。

\begin{RCode}[作业一第 1 题：drop1 偏 F 筛选]
full <- lm(glucose ~ tc + tg + insulin + hba1c, data = d)
drop1(full, test = "F")                          # 第 1 步：tc 最小，P = 0.701
m3 <- update(full, . ~ . - tc)
drop1(m3, test = "F")                            # 第 2 步：余下 3 个均 P < 0.05
anova(m3, full)                                  # 与 drop1 中 tc 一行等价
\end{RCode}

\begin{table}[htbp]
\centering\small
\caption{作业一 drop1 后退法的偏 $F$ 检验过程（$\alpha=0.05$）}
\begin{tabular}{@{}lrrrrrl@{}}
\toprule
检验变量 & 偏回归平方和 & 剔除后 SSE & $F$ & 自由度 & $P$ & 结论\\
\midrule
\multicolumn{7}{@{}l}{第 1 步：全模型，$\mathrm{SSE}=88.841$，$\mathrm{MSE}=4.038$}\\
tc      & 0.613  & 89.454  & 0.152 & $(1,22)$ & 0.701 & 剔除\\
tg      & 11.963 & 100.804 & 2.962 & $(1,22)$ & 0.099 & \\
insulin & 20.064 & 108.905 & 4.968 & $(1,22)$ & 0.036 & \\
hba1c   & 27.794 & 116.635 & 6.883 & $(1,22)$ & 0.016 & \\
\midrule
\multicolumn{7}{@{}l}{第 2 步：剔除 tc 后，$\mathrm{SSE}=89.454$，$\mathrm{MSE}=3.889$}\\
tg      & 26.530 & 115.984 & 6.821 & $(1,23)$ & 0.016 & 保留\\
insulin & 25.690 & 115.144 & 6.605 & $(1,23)$ & 0.017 & 保留\\
hba1c   & 32.269 & 121.723 & 8.297 & $(1,23)$ & 0.008 & 保留\\
\bottomrule
\end{tabular}
\end{table}

\paragraph{2. 筛选结论。}第 1 步 tc 的 $F(1,22)=0.152,P=0.701$ 最小，剔除（此时 tg 虽为 $P=0.099$，按单步剔除规则留待下一步）。第 2 步删去 tc 后，tg 的 $F$ 值从 2.962 跃升至 6.821（因为 tc 与 tg 高度相关，此前分走了 tg 的部分解释力），余下三个变量均 $P<0.05$，停止。最终模型为 $\text{glucose}\sim\text{tg}+\text{insulin}+\text{hba1c}$。

\subsubsection*{第 2 题\quad 最终模型的偏回归系数与拟合优度}
\emph{题目要求：计算最终模型各变量的偏回归系数并解释统计学意义，给出 $R^2$ 与调整 $R^2$。}

\begin{RCode}[作业一第 2 题：拟合与系数区间]
fit <- m3
summary(fit); confint(fit)
c(R2 = summary(fit)$r.squared, adjR2 = summary(fit)$adj.r.squared, s = summary(fit)$sigma)
\end{RCode}

\begin{table}[htbp]
\centering\small
\caption{作业一最终模型的偏回归系数估计（$n=27$）}
\begin{tabular}{@{}lrrrrl@{}}
\toprule
变量 & $b$ & SE & $t$ & $P$ & 95\% CI\\
\midrule
截距                  & 6.500  & 2.396 & 2.713  & 0.012 & $(1.543,\ 11.456)$\\
tg（每 1 mmol/L）     & 0.402  & 0.154 & 2.612  & 0.016 & $(0.084,\ 0.721)$\\
insulin（每 1 mU/L）  & $-0.287$ & 0.112 & $-2.570$ & 0.017 & $(-0.518,\ -0.056)$\\
hba1c（每 1\%）       & 0.663  & 0.230 & 2.880  & 0.008 & $(0.187,\ 1.140)$\\
\bottomrule
\end{tabular}
\TableNote{模型 $R^2=0.598$，调整 $R^2=0.546$，$s=1.972$ mmol/L，整体 $F(3,23)=11.41,P<0.001$。}
\end{table}

控制其余两个变量不变时：tg 每升 1 mmol/L，空腹血糖平均升 0.402 mmol/L；insulin 每升 1 mU/L，空腹血糖平均降 0.287 mmol/L；hba1c 每升 1 个百分点，空腹血糖平均升 0.663 mmol/L；三项 $t$ 检验均显著。截距 6.500 对应三变量全为 0（远在样本外），无实际意义。相比全模型，删去 tc 后 $R^2$ 几乎不变（$0.601\to 0.598$），调整 $R^2$ 反由 0.528 升至 0.546，$s$ 降至 1.972 mmol/L。需注意筛选后的 $P$ 值偏乐观，且 25 号点在最终模型中杠杆值达 0.476、Cook 距离 0.694（删去后 tg 系数降至 0.178，$P=0.362$）。

\subsubsection*{第 3 题\quad 标准化偏回归系数}
\emph{题目要求：计算标准化偏回归系数，比较相对贡献并说明与原始系数排序为何不同。}

\begin{RCode}[作业一第 3 题：标准化偏回归系数]
b <- coef(fit)[-1]; sx <- sapply(d[, names(b)], sd); sy <- sd(d$glucose)
round(b * sx / sy, 3)                            # b* = b * SD(x) / SD(y)
dz <- as.data.frame(scale(d[, -1]))              # 验证：全部标准化后再回归
coef(lm(glucose ~ tg + insulin + hba1c, data = dz))
\end{RCode}

\begin{table}[htbp]
\centering\small
\caption{作业一原始与标准化偏回归系数对照}
\begin{tabular}{@{}lrrrlcc@{}}
\toprule
变量 & $b$ & $S_x$ & $b^*$ & $b^*$ 的 95\% CI & $|b|$ 排序 & $|b^*|$ 排序\\
\midrule
tg      & 0.402  & 2.575 & 0.354  & $(0.074,\ 0.635)$   & 2 & 3\\
insulin & $-0.287$ & 3.671 & $-0.360$ & $(-0.650,\ -0.070)$ & 3 & 2\\
hba1c   & 0.663  & 1.823 & 0.413  & $(0.116,\ 0.710)$   & 1 & 1\\
\bottomrule
\end{tabular}
\end{table}

按 $|b^*|$，hba1c 最大（0.413），insulin（$-0.360$）略高于 tg（0.354），与原始系数 $|b|$ 的排序（hba1c $>$ tg $>$ insulin）不同。原因是原始系数带量纲，insulin 的标准差（3.67）大于 tg（2.57），换算到“每变化 1 个标准差”的共同尺度后两者易位。但 insulin 与 tg 的 $|b^*|$ 仅差 0.006 且置信区间高度重叠，三者处于同一量级。

\subsubsection*{第 4 题\quad 自变量取均值时的区间估计}
\emph{题目要求：计算各自变量取样本均数时均数的 95\% CI 与新个体的 95\% 预测区间。}

\begin{RCode}[作业一第 4 题：均值 CI 与个体 PI]
new <- data.frame(tg = mean(d$tg), insulin = mean(d$insulin), hba1c = mean(d$hba1c))
predict(fit, new, interval = "confidence")
predict(fit, new, interval = "prediction")
\end{RCode}

回归面必过样本中心，故 $\hat y_0=\bar y=11.926$ mmol/L，此处 $\mathrm{SE}(\hat y_0)=s/\sqrt{n}=1.972/\sqrt{27}=0.380$（所有点中最小），$t_{0.975,23}=2.069$：
\begin{itemize}
\item \textbf{均数 95\% CI}：$11.926\pm 2.069\times 0.380=(11.14,\ 12.71)$ mmol/L（宽 1.57）；
\item \textbf{个体 95\% PI}：$11.926\pm 2.069\times\sqrt{1.972^2+0.380^2}=(7.77,\ 16.08)$ mmol/L（宽 8.31，为 CI 的 5.3 倍，因为额外包含了个体随机变异 $s^2=3.889$）。
\end{itemize}

\subsubsection*{思考一与思考二\quad 总胆固醇为何不显著？逐步选择有何不足？}

\begin{RCode}[作业一思考题：偏相关、25 号影响与 Bootstrap 稳定性]
e.y  <- resid(lm(glucose ~ tg + insulin + hba1c, data = d))
e.tc <- resid(lm(tc ~ tg + insulin + hba1c, data = d))
cor(e.y, e.tc)                                   # 偏相关 = 0.083
summary(lm(glucose ~ tc, data = d[-25, ]))$coefficients
m0 <- lm(glucose ~ 1, data = d); sc <- ~ tc + tg + insulin + hba1c
add1(m0, sc, test = "F")                         # 前进法：hba1c -> tc，止于 tc + hba1c
step(m0, scope = sc, direction = "both")         # AIC 双向逐步：tg + insulin + hba1c
\end{RCode}

\begin{itemize}
\item \textbf{tc 为何从 $P=0.002$ 变为 $P=0.701$？} 一是信息重叠：tc 对其余三变量回归的 $R^2=0.542$，扣除后偏相关仅 0.083，系数从 1.026 跌至 0.142（标准误仅由 0.305 增至 0.366，并非主要因 VIF 膨胀）；二是 25 号点杠杆高（单变量 $h=0.534$），删去 25 号后连 tc 的简单回归也不再显著（$b=0.616,P=0.166$）。不能据此说 tc“不重要”：其 95\% CI 宽至 $(-0.62,0.90)$，且临床血脂控制并不取决于对空腹血糖的额外预测力。
\item \textbf{逐步选择的四点不足：} (1) 依赖方法与阈值——前进法（$\alpha_{\text{入}}=0.05$）停在 \texttt{tc + hba1c}（调整 $R^2=0.441$），与后退法完全不同；(2) 结果不稳定——1000 次 bootstrap 后退法中，\texttt{tg + insulin + hba1c} 仅占 27.4\%，各变量入选率依次为 hba1c 70.1\%、insulin 64.3\%、tg 53.5\%、tc 27.0\%；(3) 筛选后 $P$ 值偏乐观、区间偏窄；(4) 忽视专业混杂结构。改进办法：按专业知识预设混杂变量、小变量数用全子集与交叉验证、高相关时改用惩罚回归（岭回归/LASSO）并报告 bootstrap 入选率。
\end{itemize}

% ============================================================================
\subsection{作业二：全模型的回归诊断与 datasetB 的岭回归}

\subsubsection*{第 1 题\quad diabetes27 全模型的回归诊断}
\emph{题目要求：对 \texttt{diabetes27\_hw.csv}（$n=27$）的全模型 $\widehat{\text{glucose}}=5.943+0.142\,\text{tc}+0.351\,\text{tg}-0.271\,\text{insulin}+0.638\,\text{hba1c}$ 逐项检查 LINE 四条件、共线性、强影响点及删除前后变化，并讨论真实高脂血症患者是否应删除。}

\paragraph{1. LINE 四条件。}
\begin{RCode}[作业二第 1 题：LINE 诊断与二次项检验]
library(car); library(MASS)
d1 <- read.csv("data/diabetes27_hw.csv")
m <- lm(glucose ~ tc + tg + insulin + hba1c, data = d1)
par(mfrow = c(2, 2)); plot(m)
shapiro.test(resid(m)); ncvTest(m)
set.seed(2026); durbinWatsonTest(m, reps = 5000)
residualPlots(m, plot = FALSE)                   # Tukey 曲率检验
anova(m, update(m, . ~ . + I(hba1c^2)))          # 补查 hba1c 二次项
\end{RCode}

\begin{figure}[htbp]
\centering
\includegraphics[width=0.92\textwidth]{figures/hw2_diag.pdf}
\caption{作业二第 1 题全模型的四张诊断图（蓝线为平滑线，橙色虚线为 Cook 距离等高线）}
\end{figure}

\begin{table}[htbp]
\centering\small
\caption{作业二第 1 题 LINE 四条件的检查结果}
\begin{tabular}{@{}lp{6.4cm}p{4.8cm}@{}}
\toprule
条件 & 依据 & 结论\\
\midrule
L 线性 & 残差--拟合值图平滑线呈 U 形；Tukey 曲率检验统计量 $2.89,P=0.004$ & 存疑，加入 $\text{hba1c}^2$ 后 $F(1,21)=6.85,P=0.016$（删 25 号仍 $P=0.011$）\\
I 独立 & $\mathrm{DW}=1.635,P=0.353$；但横断面数据行序无意义 & DW 不作依据；27 例为不同患者，无聚类/重复，按独立处理\\
N 正态 & Q--Q 图点基本贴近对角线；Shapiro--Wilk $W=0.978,P=0.823$ & 满足\\
E 等方差 & 尺度--位置图无喇叭形；ncvTest $\chi^2=0.0004,P=0.984$ & 未见异方差（$n=27$ 效能有限）\\
\bottomrule
\end{tabular}
\end{table}

\paragraph{2. 共线性。}四个变量的 VIF 依次为 tc 2.186、tg 1.780、insulin 1.278、hba1c 1.267，最大相关系数 0.632（tc 与 tg），条件数 $\kappa=2.67$，无严重共线性。tc 进入全模型后 $P$ 值从 0.002 升至 0.701，主要不是标准误膨胀（标准误仅放大 1.20 倍，若系数不变仍显著），而是调整其他变量后系数本身缩水 86\%（偏相关 0.083）以及 25 号点的高杠杆所致。

\paragraph{3. 强影响点与删除对照。}
\begin{RCode}[作业二第 1 题：杠杆值、删除学生化残差与 Cook 距离]
inf <- data.frame(id = d1$id, h = hatvalues(m), r = rstudent(m), D = cooks.distance(m))
inf[order(-inf$D)[1:4], ]; outlierTest(m)
m.del <- update(m, data = d1[-25, ])             # 删除 25 号对比
m.mm  <- rlm(glucose ~ tc + tg + insulin + hba1c, data = d1, method = "MM")
\end{RCode}

\begin{table}[htbp]
\centering\small
\caption{作业二第 1 题强影响点（上）与删除 25 号前后的系数对照（下）}
\begin{tabular}{@{}ccccccccc@{}}
\toprule
id & tc & tg & insulin & hba1c & glucose & $h_i$ & $r_i^*$ & $D_i$\\
\midrule
25 & 11.54 & 10.89 & 1.20 & 10.5 & 20.0 & 0.613 & 1.864  & 0.990\\
6  & 6.05  & 0.64  & 1.42 & 13.6 & 18.3 & 0.359 & 1.965  & 0.383\\
13 & 7.97  & 1.93  & 7.57 & 9.9  & 8.4  & 0.274 & $-2.319$ & 0.338\\
26 & 5.84  & 0.92  & 8.61 & 6.4  & 13.3 & 0.203 & 2.853  & 0.314\\
\bottomrule
\end{tabular}

\vspace{0.8em}
\begin{tabular}{@{}lrrrrl@{}}
\toprule
变量 & 删除前 $b$ & 删除前 $P$ & 删除后 $b$ & 删除后 $P$ & 变化\\
\midrule
tc      & 0.142  & 0.701 & $-0.243$ & 0.554 & 符号翻转\\
tg      & 0.351  & 0.099 & 0.220  & 0.298 & 减小 37\%\\
insulin & $-0.271$ & 0.036 & $-0.234$ & 0.058 & 绝对值降 13\%，$P$ 越过 0.05\\
hba1c   & 0.638  & 0.016 & 0.752  & 0.005 & 增大 18\%\\
\midrule
$R^2$   & 0.601  &       & 0.508  &       & 下降 0.093\\
\bottomrule
\end{tabular}
\TableNote{界值：$h_i>0.370$，$|r_i^*|>2$，$D_i>0.182$。25 号是高杠杆点而非离群点；删除学生化残差最大的是 26 号（$r^*=2.853$，Bonferroni 校正 $P=0.257$）。}
\end{table}

\paragraph{4. 思考：真实高脂血症患者是否删除？}不删。25 号是真实存在的重度高脂血症患者，删除会人为缩小目标人群。规范做法是主分析保留该例、并排报告删除前后对照表，明确说明 hba1c 结论稳健而 tc 方向、tg 幅度与 insulin 显著性依赖 25 号，同时注明其余 26 人 tg 最高 8.50、高值区不可外推。MM 稳健回归给 25 号的权重高达 0.89（因残差不大），无法自动降权，仍须依赖敏感性分析。

\subsubsection*{第 2 题\quad datasetB 的岭回归}
\emph{题目要求：对 \texttt{datasetB.csv}（$n=150$）找出共线性变量，按“系数合理且岭迹平稳”选出岭参数并建方程，同时评估强影响点效应。}

\paragraph{1. 共线性与强影响点。}OLS 拟合得 $\widehat{\text{glu}}=2.462+0.0314\,\text{age}+0.0671\,\text{bmi}-0.0063\,\text{waist}$（$R^2=0.459$）。bmi 与 waist 的相关系数高达 $r=0.968$，VIF 分别为 16.02 与 16.14（age 为 1.03），条件数 $\kappa=7.93$，导致二者在模型中均不显著且 waist 符号反常为负（单独与 age 配对时二者均 $P<0.001$）。诊断出两个强影响点：149 号（$h=0.297,r^*=2.938,D=0.864$）恰恰偏离了 bmi--waist 共线方向（比其余 149 人的 waist $\sim$ bmi 直线高出 10.4 cm）；150 号（$h=0.032,r^*=6.847,D=0.292$）则是纯 $y$ 方向离群点。

\paragraph{2. 岭参数选择与敏感性分析。}
\begin{RCode}[作业二第 2 题：四种样本岭迹与敏感性比较]
d2 <- read.csv("data/datasetB.csv")
lam <- seq(0, 100, 0.5)
sets <- list(all = 1:150, del149 = -149, del150 = -150, del2 = -(149:150))
rd <- lapply(sets, function(i) lm.ridge(glu ~ age + bmi + waist, d2[i, ], lambda = lam))
sapply(rd, function(r) lam[which.min(r$GCV)])    # GCV 最优 lambda
rid <- lm.ridge(glu ~ age + bmi + waist, d2, lambda = 30)
coef(rid)                                        # 原始尺度最终方程
\end{RCode}

\begin{figure}[htbp]
\centering
\includegraphics[width=0.92\textwidth]{figures/hw2_ridge.pdf}
\caption{作业二第 2 题四种样本下的岭迹图（纵轴为自变量标准化后的系数，圆圈标示 waist 转正处）}
\end{figure}

\begin{table}[htbp]
\centering\small
\caption{作业二第 2 题岭参数下限（上）与 OLS / 岭回归敏感性对照（下）}
\begin{tabular}{@{}lccc@{}}
\toprule
样本 & GCV 最小 $\lambda$ & 系数合理（bmi、waist 均正） & 岭迹平稳（$\Delta\lambda=10$ 变化 $\le 0.01$）\\
\midrule
全部 150 人      & 7.5 & 3.0  & 18.5\\
删 149 号        & 1.5 & 18.5 & 30.0\\
删 150 号        & 6.0 & 3.0  & 20.0\\
删 149、150 号   & 1.0 & 15.5 & 30.5\\
\bottomrule
\end{tabular}

\vspace{0.8em}
\begin{tabular}{@{}lrrrrrr@{}}
\toprule
模型与样本 & 截距 & age & bmi & waist & $R^2$ & bmi+waist（标化）\\
\midrule
OLS · 全部 150 人     & 2.4617 & 0.0314 & 0.0671 & $-0.0063$ & 0.4594 & 0.166\\
OLS · 删 149 号       & 3.4314 & 0.0336 & 0.1241 & $-0.0369$ & 0.4791 & 0.110\\
OLS · 删 149、150 号  & 3.3365 & 0.0307 & 0.1263 & $-0.0349$ & 0.5154 & 0.131\\
\midrule
岭（$\lambda=30$）· 全部 & 2.4546 & 0.0260 & 0.0283 & 0.0083 & 0.4469 & 0.158\\
岭（$\lambda=30$）· 删 149 & 2.7785 & 0.0270 & 0.0297 & 0.0031 & 0.4575 & 0.113\\
岭（$\lambda=30$）· 删两人 & 2.6961 & 0.0246 & 0.0329 & 0.0043 & 0.4914 & 0.132\\
\bottomrule
\end{tabular}
\end{table}

\paragraph{3. 结论与局限。}综合四种样本取 $\lambda=30$（对应相关阵 $k=\lambda/n=0.20$；GCV 最小 7.5 处 waist 正号依赖 149 号，故不采纳），全样本原始尺度回归方程为
\[
\widehat{\text{glu}}=2.455+0.0260\,\text{age}+0.0283\,\text{bmi}+0.0083\,\text{waist}
\]
（$R^2=0.447$，留一交叉验证 MSE 0.311 略优于 OLS 的 0.315）。$\lambda=30$ 显著压制了 149 号的影响（杠杆值 $0.297\to 0.170$，类 Cook 距离 $0.864\to 0.272$），但管不了 $y$ 方向离群的 150 号（类 Cook 距离 $0.292\to 0.313$），且对 age 造成了 17.3\% 的收缩偏差；岭回归也不解决数据本身的异方差（ncvTest $P<0.001$）。

% ============================================================================
\subsection{作业三：3×3 析因设计方差分析}

数据为 \texttt{sbp\_factorial3x3.csv}（$N=54$，两因素药物 A 剂量 10/20/30 mg $\times$ 药物 B 剂量 2.5/5/10 mg，共 9 组，每组 $n=6$，结局为 SBP 下降值，单位 mmHg）。

\begin{RCode}[作业三：读入、效应编码与前提检查]
library(car)
options(contrasts = c("contr.sum", "contr.poly"))
d <- read.csv("data/sbp_factorial3x3.csv", check.names = FALSE, fileEncoding = "UTF-8")
names(d) <- c("id", "A", "B", "y")
d$A <- factor(substr(d$A, 1, 2)); d$B <- factor(substr(d$B, 1, 2))
m <- lm(y ~ A * B, data = d)
leveneTest(y ~ A * B, data = d); shapiro.test(resid(m)); outlierTest(m)
\end{RCode}

\subsubsection*{第 1 题\quad 设计类型、统计模型与前提检查}
\emph{题目要求：指出设计类型、因素水平与试验单元；写出固定效应模型、符号含义与基本假定，并检查假定。}

本试验为完全随机设计的平衡 $3\times 3$ 析因设计（9 个处理组，每组 6 例，54 名患者即 54 个独立试验单元）。固定效应模型为
\[
y_{ijk}=\mu+\alpha_i+\beta_j+(\alpha\beta)_{ij}+\varepsilon_{ijk},\qquad i,j=1,2,3;\ k=1,\dots,6,
\]
约束 $\sum_i\alpha_i=\sum_j\beta_j=0$、$\sum_i(\alpha\beta)_{ij}=\sum_j(\alpha\beta)_{ij}=0$，共 9 个自由参数，假定 $\varepsilon_{ijk}\overset{\text{iid}}{\sim}N(0,\sigma^2)$。

\begin{figure}[htbp]
\centering
\includegraphics[width=0.92\textwidth]{figures/hw3_resid.pdf}
\caption{作业三各处理组残差分布与学生化残差正态 Q--Q 图（橙色点为 28 号）}
\end{figure}

前提检查：方差齐性 Levene 检验 $F(8,45)=0.85,P=0.567$；正态性 Shapiro--Wilk $W=0.952,P=0.030$，问题仅出在 28 号（A2B2 组，$y=4.96$，组均 14.00，删除学生化残差 $-4.00$，Bonferroni $P=0.013$）。删去 28 号后用 II 型平方和重算，$W=0.958$（$P=0.057$），A、B 主效应仍 $P<0.001$，交互 $F(4,44)=4.85,P=0.002$，结论稳健，故主分析保留全部 54 例。

\subsubsection*{第 2 题\quad 假设检验与方差分析表}
\emph{题目要求：写出三个原假设并完成方差分析表。}

原假设依次为 $H_0^A:\alpha_1=\alpha_2=\alpha_3=0$、$H_0^B:\beta_1=\beta_2=\beta_3=0$、$H_0^{AB}:\text{全部 }(\alpha\beta)_{ij}=0$。平衡设计下 I、II、III 型平方和一致。

\begin{RCode}[作业三第 2 题：方差分析与手算核对]
anova(m); Anova(m, type = 3)
cm <- tapply(d$y, list(d$A, d$B), mean); cs <- tapply(d$y, list(d$A, d$B), sd)
gm <- mean(cm); am <- rowMeans(cm); bm <- colMeans(cm)
c(SSA = 18 * sum((am - gm)^2), SSB = 18 * sum((bm - gm)^2), SSE = 5 * sum(cs^2))
\end{RCode}

\begin{table}[htbp]
\centering\small
\caption{作业三 $3\times 3$ 析因方差分析表}
\begin{tabular}{@{}lrrrrll@{}}
\toprule
变异来源 & SS & df & MS & $F$ & $P$ & 显著性\\
\midrule
因素 A          & 1225.57 & 2  & 612.78 & 74.92 & $<0.001$ & 显著\\
因素 B          & 789.41  & 2  & 394.71 & 48.26 & $<0.001$ & 显著\\
A$\times$B      & 106.78  & 4  & 26.70  & 3.26  & $0.020$  & 显著\\
误差            & 368.06  & 45 & 8.18   &       &          & \\
\midrule
总计            & 2489.82 & 53 &        &       &          & \\
\bottomrule
\end{tabular}
\TableNote{模型 $R^2=0.852$；A 占总变异 49.2\%，B 占 31.7\%，A$\times$B 占 4.3\%，交互作用显著但量级远小于两主效应。}
\end{table}

\subsubsection*{第 3 题与第 4 题\quad 主效应趋势、交互模式与单独效应}
\emph{题目要求：结合边际均数解释主效应；结合 9 格均数剖析交互模式、专业含义及交互显著时如何解读主效应。}

\begin{RCode}[作业三第 3–4 题：Tukey 比较、正交多项式趋势与单独效应]
TukeyHSD(aov(y ~ A * B, data = d), c("A", "B"))
mse <- deviance(m) / df.residual(m)
contr1 <- function(w) {
  L <- sum(w * cm); SS <- 6 * L^2 / sum(w^2)
  c(L = L, SE = sqrt(mse * sum(w^2) / 6), SS = SS, F = SS / mse,
    P = pf(SS / mse, 1, 45, lower.tail = FALSE))
}
lin <- c(-1, 0, 1); quad <- c(1, -2, 1); one <- rep(1, 3)
contr1(outer(lin, lin))                          # 线性 x 线性交互对比
\end{RCode}

\begin{table}[htbp]
\centering\small
\caption{作业三单元格均数（括号内为交互效应）、增量与单独效应（mmHg）}
\begin{tabular}{@{}lccccc@{}}
\toprule
 & B 2.5 mg & B 5 mg & B 10 mg & B 10 $-$ 2.5 mg & B 的单独效应 $F$\\
\midrule
A 10 mg & 4.00 ($+1.56$)  & 8.00 ($+0.22$)  & 10.00 ($-1.78$) & 6.00  & 6.85\\
A 20 mg & 9.00 ($+0.56$)  & 14.00 ($+0.22$) & 17.00 ($-0.78$) & 8.00  & 11.98\\
A 30 mg & 12.00 ($-2.11$) & 19.00 ($-0.45$) & 26.00 ($+2.56$) & 14.00 & 35.96\\
\midrule
A 30 $-$ 10 mg   & 8.00  & 11.00 & 16.00 & & \\
A 的单独效应 $F$ & 11.98 & 22.25 & 47.21 & & \\
\bottomrule
\end{tabular}
\TableNote{边际均数：A 依次 7.33、13.33、19.00；B 依次 8.33、13.67、17.67（SE 均为 0.67）。单独效应 $F(2,45)$ 全部显著（$P\le 0.003$）。}
\end{table}

\begin{figure}[htbp]
\centering
\includegraphics[width=0.92\textwidth]{figures/hw3_inter.pdf}
\caption{作业三 9 个处理组的均数轮廓图（误差线为 $\pm 1$ SE；左图横轴为对数剂量刻度）}
\end{figure}

\begin{itemize}
\item \textbf{主效应量效关系}：A 每增 10 mg 平均多降 5.8 mmHg（线性项占 $SS_A$ 的 99.97\%，二次项 $P=0.84$）；B 剂量每翻一倍平均多降 4.7 mmHg（$\log_2$ 线性项占 $SS_B$ 的 99.3\%，二次项 $P=0.42$）。因无单药组和安慰剂组，主效应仅代表联合用药下的加量效果。
\item \textbf{有序协同交互}：三条折线同向单调上升、互不交叉，但向右逐渐张开——属于\textbf{有序（量的）交互作用}。交互 4 个自由度中，线性$\times$线性对比（$14.0-6.0=8.00$ mmHg，95\% CI 3.30--12.71，$F(1,45)=11.75,P=0.001$）独占 $SS_{AB}$ 的 90.0\%，其余 3 个自由度合并 $P=0.73$。临床上 A3B3 降压最大（26.0 mmHg，比可加预测高 2.6 mmHg），但需警惕低血压及剂量相关不良反应。
\item \textbf{交互显著时如何理解主效应}：因为是有序交互且 6 个单独效应同向显著，主效应的方向依然可信（加量总能增强降压），但具体幅度必须按另一因素水平分列报告；若换作交叉的无序交互，则不能孤立解读主效应。
\end{itemize}

% ============================================================================
\subsection{作业四：三种康复方案的协方差分析与脑卒中重复测量}

\subsubsection*{第 1 题\quad 三种康复方案的协方差分析}
\emph{题目要求：对 \texttt{knee3.csv}（$N=90$，A 常规、B 强化、C 机器人辅助各 30 例），以术前评分为协变量、术后 6 个月评分为结局，依次做斜率同质性检验、整体 $F$ 检验和 Bonferroni 两两比较。}

\paragraph{1. 数据描述与斜率同质性。}三组术前评分并不均衡（A 57.94、B 64.06、C 61.42，$F(2,87)=3.65,P=0.030$），且术前与术后高度相关（各组 $r=0.77,0.84,0.87$），直接比术后评分会把 B 组起点高的优势误算为疗效。拟合含交互项的完整模型检验斜率同质性：

\begin{RCode}[作业四第 1 题：协方差分析完整流程]
library(car); library(emmeans)
k <- read.csv("data/knee3.csv", fileEncoding = "UTF-8")[, -1]
k$group <- factor(k$group, labels = c("A", "B", "C"))
anova(lm(post ~ pre * group, data = k))          # 1. 斜率同质性：P = 0.758
fit <- lm(post ~ pre + group, data = k)          # 2. 平行模型
Anova(fit, type = 2)                             #    II 型平方和，与顺序无关
emm <- emmeans(fit, ~ group)                     # 3. 在 pre = 61.14 处的校正均数
pairs(emm, adjust = "bonferroni"); confint(pairs(emm, adjust = "bonferroni"))
\end{RCode}

交互项 $F(2,84)=0.28,P=0.758$，三组各自斜率分别为 1.000、1.074、0.947，满足平行假定；其余前提亦通过（二次项 $P=0.385$，Shapiro--Wilk $P=0.656$，Levene $P=0.795$，最大 Cook 距离 0.078）。

\paragraph{2. 整体检验与校正均数比较。}
\begin{table}[htbp]
\centering\small
\caption{作业四第 1 题协方差分析表（上）及校正前后均数与两两比较（下）}
\begin{tabular}{@{}lrrrrl@{}}
\toprule
变异来源（II 型 SS） & SS & df & MS & $F$ & $P$\\
\midrule
术前评分（协变量） & 6767.9 & 1  & 6767.9 & 190.59 & $<0.001$\\
康复方案           & 570.4  & 2  & 285.2  & 8.03   & $0.0006$\\
残差               & 3053.9 & 86 & 35.5   &        & \\
\bottomrule
\end{tabular}

\vspace{0.8em}
\begin{tabular}{@{}lrrrrl@{}}
\toprule
比较 & 未校正差值 & 未校正 $P$ & 校正后差值 & Bonferroni 95\% CI & 校正后 $P$\\
\midrule
C $-$ A & 7.53 & 0.022 & 4.04  & $(0.23,\ 7.85)$   & 0.034\\
C $-$ B & 3.43 & 0.642 & 6.07  & $(2.29,\ 9.86)$   & 0.0005\\
B $-$ A & 4.10 & 0.416 & $-2.03$ & $(-5.94,\ 1.88)$ & 0.625\\
\bottomrule
\end{tabular}
\TableNote{共同斜率 $b=1.002$，$R^2=0.714$，方案偏 $\eta^2=0.157$。在 $\bar x=61.14$ 处的校正均数（SE）：A 78.64 (1.11)、B 76.61 (1.11)、C 82.68 (1.09)。误差均方较单因素方差分析（112.9）压缩了 69\%。}
\end{table}

\begin{figure}[htbp]
\centering
\includegraphics[width=0.92\textwidth]{figures/hw4_ancova.pdf}
\caption{作业四第 1 题：(a) 三组术前--术后散点及回归直线；(b) 校正前后均数及 95\% CI}
\end{figure}

\paragraph{3. 结论与提示。}机器人辅助康复（C）效果最好，校正后分别比 A 高 4.0 分、比 B 高 6.1 分；B 组未校正时看似比 A 高 4.1 分，实为术前起点高 6.1 分所致，扣除后反而略低 2.0 分（无显著差异）。未校正分析漏掉了 C 优于 B 的差别。用 \texttt{anova()} 时务必把协变量 \texttt{pre} 放在 \texttt{group} 前面，或直接用 \texttt{car::Anova(..., type = 2)}。

\subsubsection*{第 2 题\quad 脑卒中患者上肢功能评分的重复测量方差分析}
\emph{题目要求：随机对照试验 \texttt{stroke.csv}（$N=60$，新型训练组与常规训练组各 30 例，在入组 T0、4 周 T4、8 周 T8 测上肢功能评分），按“画图 → 验球形 → 看交互 → 比两两”四步完成分析。}

\begin{RCode}[作业四第 2 题：转长表、重复测量 ANOVA 与简单效应]
library(afex); library(tidyr); library(emmeans)
s <- read.csv("data/stroke.csv", fileEncoding = "UTF-8")
s$group <- factor(s$group, levels = c("新型训练组", "常规训练组"))
L <- pivot_longer(s, cols = c(baseline, y_wk4, y_wk8), names_to = "time", values_to = "y")
L$time <- factor(L$time, levels = c("baseline", "y_wk4", "y_wk8"), labels = c("T0", "T4", "T8"))
mod <- aov_ez(id = "id", dv = "y", data = L, within = "time", between = "group", type = 3)
summary(mod)                                     # Mauchly 检验与 GG/HF 校正
pairs(emmeans(mod, ~ time | group), adjust = "bonferroni")   # 组内
pairs(emmeans(mod, ~ group | time))                          # 组间
contrast(emmeans(mod, ~ time * group), interaction = c("consec", "pairwise"))
\end{RCode}

\paragraph{1. 描述与画图。}两组基线可比（入组评分 44.60 vs 46.20，$P=0.277$），8 周内新型组由 44.60 升至 56.10（4 周）、63.53（8 周），净增 18.93 分；常规组由 46.20 升至 52.47、56.80，净增 10.60 分。两折线越往后张得越开（提示交互），且 8 周时标准差明显大于入组（提示方差齐性/球形可能不成立）。

\begin{figure}[htbp]
\centering
\includegraphics[width=0.95\textwidth]{figures/hw4_profile.pdf}
\caption{作业四第 2 题：(a) 两组各时点均数与 95\% CI；(b)(c) 60 名患者个体轨迹与组均数}
\end{figure}

\paragraph{2. 球形检验与三个 $F$ 检验。}两组合并协方差阵的 Box's $M$ 检验 $P=0.140$（可合并）。三个差值方差分别为 20.33（T0--T4）、37.98（T0--T8）与 58.21（T4--T8），相差近 3 倍；Mauchly 检验 $W=0.683,\chi^2(2)=21.77,P=1.9\times 10^{-5}$，拒绝球形假设，采用 Greenhouse--Geisser 校正（$\hat\varepsilon_{\mathrm{GG}}=0.759$，HF 为 0.775）。

\begin{table}[htbp]
\centering\small
\caption{作业四第 2 题重复测量方差分析表（上）与简单效应及交互对比（下）}
\begin{tabular}{@{}lrrrrll@{}}
\toprule
变异来源 & SS & df & MS & $F$ & 未校正 $P$ & GG 校正 $P$\\
\midrule
训练方案（个体间） & 384.27   & 1   & 384.27  & 3.51   & 0.066    & ---\\
受试者间误差       & 6351.61  & 58  & 109.51  &        &          & \\
时间（个体内）     & 6631.63  & 2   & 3315.82 & 170.74 & $<0.001$ & $<0.001$\\
方案 $\times$ 时间 & 532.21   & 2   & 266.11  & 13.70  & $<0.001$ & $<0.001$\\
受试者内误差       & 2252.82  & 116 & 19.42   &        &          & \\
\bottomrule
\end{tabular}

\vspace{0.8em}
\begin{tabular}{@{}lllrl@{}}
\toprule
分析方向 & 比较 & 差值 & 95\% CI & $P$\\
\midrule
组内（新型组） & T4 $-$ T0 / T8 $-$ T0 / T8 $-$ T4 & 11.50 / 18.93 / 7.43 & 均不含 0 & 均 $<0.001$\\
组内（常规组） & T4 $-$ T0 / T8 $-$ T0 / T8 $-$ T4 & 6.27 / 10.60 / 4.33  & 均不含 0 & $<0.001$ / $<0.001$ / $0.009$\\
组间（新$-$常）& T0 / T4 / T8                      & $-1.60$ / $3.63$ / $6.73$ & $(-4.52,1.32)$ / $(0.28,6.99)$ / $(2.28,11.19)$ & $0.277$ / $0.034$ / $0.004$\\
交互对比       & 0--4 周 / 4--8 周 / 0--8 周       & 5.23 / 3.10 / 8.33   & $(2.90,7.56)$ / $(-0.84,7.04)$ / $(5.15,11.52)$ & $<0.001$ / $0.121$ / $<0.001$\\
\bottomrule
\end{tabular}
\TableNote{GG 校正后时间与交互自由度均为 $(1.52,88.05)$，交互精确 $P=4.2\times 10^{-5}$（多元 Pillai 迹 $P<0.001$，结论一致）。若对三个时点的组间比较合并做 Bonferroni 校正，T4 为 $P=0.103$、T8 仍 $P=0.011$。}
\end{table}

\paragraph{3. 结论与解读。}方案 $\times$ 时间交互显著（偏 $\eta^2=0.19$），说明新型训练组恢复明显快于常规组：8 周净改善多出 8.3 分（95\% CI 5.2--11.5），8 周时组间高出 6.7 分（95\% CI 2.3--11.2），差距主要在前 4 周形成（多改善 5.2 分，$P<0.001$）。方案主效应不显著（$P=0.066$）只是因为把入组无差别的基线也平均进了三个时点，切不可据此误判两方案疗效相同。

\begin{notebox}[四份作业代码的共同写作模板]
回头看四份完整脚本，其实遵循同一套可复用的分析流水线：读数据并用 \Verb|str()| 核对 $\to$ 描述统计与题目表核对 $\to$ 主模型拟合与假设检验 $\to$ 前提条件检查 $\to$ 强影响点/异常值敏感性分析 $\to$ 矢量图输出。把这套模板记熟，做任何线性模型实战都不会漏步。
\end{notebox}

\appendix
\section{函数速查表}
\label{app:cheat}
按“要做的事”查函数。包名写在函数前，\Verb|stats| 与 \Verb|graphics| 是 R 自带的，不用 \Verb|library()|。

\begingroup\footnotesize
\begin{longtable}{@{}p{3.3cm}p{6.6cm}p{4.6cm}@{}}
\caption{任务 → R 函数 → 读哪个数}\\
\toprule
任务 & 写法 & 读什么\\
\midrule
\endfirsthead
\toprule
任务 & 写法 & 读什么\\
\midrule
\endhead
\bottomrule
\endfoot
\multicolumn{3}{@{}l}{\textbf{拟合与输出}}\\
拟合线性模型 & \Verb|fit <- lm(y ~ x1 + x2, data = d)| & ——\\
系数表、$R^2$、整体 F & \Verb|summary(fit)| & Estimate、SE、P；R²；F 行\\
系数 CI & \Verb|confint(fit, level = 0.95)| & 2.5\%、97.5\%\\
均数 CI / 个体 PI & \Verb|predict(fit, newdata, interval = "confidence"/"prediction")| & fit、lwr、upr\\
标准化系数 & \Verb|coef(fit)[-1] * sapply(X, sd) / sd(y)| & 同一模型内比较\\
偏相关 & \Verb|ppcor::pcor.test(x, y, z)|；或残差对残差的 \Verb|cor()| & estimate、P\\
\midrule
\multicolumn{3}{@{}l}{\textbf{检验与模型比较}}\\
嵌套模型 F 检验 & \Verb|anova(small, big)| & F、Pr(>F)\\
逐个删 / 加变量 & \Verb|drop1(fit, test = "F")|、\Verb|add1(fit, scope, test = "F")| & 每行的 F、P\\
I 型平方和 & \Verb|anova(fit)| & 注意顺序\\
II / III 型平方和 & \Verb|car::Anova(fit, type = "II")|；III 型先 \Verb|options(contrasts = c("contr.sum", "contr.poly"))| & ——\\
信息准则 & \Verb|AIC(fit)|、\Verb|BIC(fit)|、\Verb|extractAIC(fit)| & 越小越好\\
\midrule
\multicolumn{3}{@{}l}{\textbf{变量筛选}}\\
全子集 & \Verb|leaps::regsubsets(y ~ ., data, nbest = 1)| & \Verb|summary()$adjr2|、\Verb|$cp|、\Verb|$bic|\\
按 AIC / BIC 逐步 & \Verb|step(m0, scope = list(lower = ~1, upper = formula(full)), direction = "both", k = 2 或 log(n))| & 最终模型\\
按 P 值逐步 & \Verb|olsrr::ols_step_forward_p(fit, p_val = )|、\Verb|ols_step_both_p(fit, p_enter = , p_remove = )| & 入选过程\\
\midrule
\multicolumn{3}{@{}l}{\textbf{回归诊断}}\\
四张诊断图 & \Verb|par(mfrow = c(2, 2)); plot(fit)| & 形状\\
正态 & \Verb|shapiro.test(resid(fit))|、\Verb|car::qqPlot(fit)| & W、P；两端\\
等方差 & \Verb|car::ncvTest(fit)|、\Verb|car::spreadLevelPlot(fit)| & χ²、P\\
独立（有序数据） & \Verb|car::durbinWatsonTest(fit)|、\Verb|lmtest::dwtest(fit)| & DW≈2\\
线性 & \Verb|car::residualPlots(fit)|、\Verb|crPlots()|、\Verb|avPlots()| & 曲率检验 P\\
加权最小二乘 & \Verb|lm(..., weights = w)| & 再做一次诊断\\
稳健标准误 & \Verb|lmtest::coeftest(fit, vcov = sandwich::vcovHC(fit, type = "HC3"))| & 新的 SE、P\\
杠杆 / 残差 / Cook & \Verb|hatvalues()|、\Verb|rstandard()|、\Verb|rstudent()|、\Verb|cooks.distance()| & 界值见第 \ref{sec:influence} 节\\
影响点总览 & \Verb|car::influencePlot(fit)|、\Verb|car::outlierTest(fit)|、\Verb|influence.measures(fit)| & ——\\
删点重拟合 & \Verb|update(fit, data = d[-k, ])| & 系数变化\\
稳健回归 & \Verb|MASS::rlm(formula, data, method = "MM")| & 系数、权重 \Verb|$w|\\
\midrule
\multicolumn{3}{@{}l}{\textbf{共线性与岭回归}}\\
VIF & \Verb|car::vif(fit)| & > 10 严重\\
条件数 & \Verb|sqrt(max(ev)/min(ev))|，\Verb|ev <- eigen(cor(X))$values| & 15–30 以上\\
岭回归 & \Verb|MASS::lm.ridge(formula, data, lambda = seq(...))| & \Verb|select()|、\Verb|$coef|、\Verb|coef()|\\
岭迹图 & \Verb|matplot(rr$lambda, t(rr$coef), type = "l")| & 平稳的位置\\
\midrule
\multicolumn{3}{@{}l}{\textbf{一般线性模型}}\\
t 检验 & \Verb|t.test(y ~ g, data, var.equal = TRUE)| & 等价于 \Verb|lm(y ~ g)|\\
设计矩阵 & \Verb|model.matrix(fit)| & X 的列\\
换参照组 & \Verb|relevel(f, ref = "...")|、\Verb|factor(f, levels = c(...))| & 只换对比\\
单因素多重比较 & \Verb|TukeyHSD(aov(...))|、\Verb|pairwise.t.test(y, g, p.adjust.method = "bonferroni")| & 校正 P\\
交互图 & \Verb|interaction.plot(x, trace, y)| & 是否平行\\
单元格均数 & \Verb|tapply(y, list(a, b), mean)| & ——\\
简单效应 & \Verb!pairs(emmeans(fit, ~ a | b))! & 每个 b 下 a 的差\\
\midrule
\multicolumn{3}{@{}l}{\textbf{协方差分析}}\\
斜率同质性 & \Verb|anova(lm(y ~ g + x), lm(y ~ g * x))| & P > 0.05\\
校正后组间检验 & \Verb|car::Anova(lm(y ~ x + g), type = "II")| & g 一行\\
校正均数与比较 & \Verb|emm <- emmeans(fit, ~ g)|；\Verb|pairs(emm, adjust = "bonferroni")| & emmean、estimate、P\\
\midrule
\multicolumn{3}{@{}l}{\textbf{重复测量}}\\
宽转长 & \Verb|tidyr::pivot_longer(d, cols, names_to, values_to)| & ——\\
重复测量方差分析 & \Verb|afex::aov_ez(id, dv, data, within, between)|；\Verb|summary()| & F、Mauchly、GG/HF\\
校正后的表 & \Verb|anova(mod, correction = "GG", es = "pes")| & 校正自由度、P\\
基础 R 写法 & \Verb|aov(y ~ g * t + Error(id / t))| & 分层误差\\
Mauchly 检验 & \Verb|mauchly.test(lm(Y ~ 1), X = ~ 1)|（Y 为宽型矩阵） & W、P\\
两个方向的简单效应 & \Verb!emmeans(mod, ~ time | group)!、\Verb!emmeans(mod, ~ group | time)! & pairs()\\
交互对比 & \Verb|contrast(emmeans(mod, ~ time * group), interaction = c("consec", "pairwise"))| & 差距的变化\\
\end{longtable}
\endgroup

\section{课件笔记例题的复现}
\label{app:notes}
《应用线性回归\_课件笔记》各章的 R 代码块，数据原先直接写在代码里。现在数据存为 \Verb|data/| 下的 CSV，代码整理在 \Verb|scripts/09_课件笔记例题复现.R|，可以从头跑到尾。笔记修订记录中提到当时本机没装 leaps 和 ez，这两段未能直接运行；现在两个包都已安装，全部代码块均已运行，结果与笔记给出的数值一致：

\begin{table}[htbp]
\centering\small
\caption{课件笔记例题与复现结果}
\begin{tabular}{@{}p{3.2cm}p{3.8cm}p{7.0cm}@{}}
\toprule
例题 & 数据文件 & 复现的关键数值\\
\midrule
第 1 章 矩阵 & body12.csv & 均数向量、协方差阵、相关阵、特征值\\
例 3.1 肺活量 & lung29.csv & $b=(-0.5657, 0.005017, 0.05406)$；$Q=2.558$，$s^2=0.0984$\\
例 15-1 全子集与 AIC & diabetes27\_hw.csv & \Verb|which.max(adjr2)| = 11；AIC = 120.78；extractAIC = 42.157\\
偏回归图 & 同上 & 斜率 −0.2870 = 全模型中 x3 的系数\\
岭回归 & 同上 & GCV 最优 λ = 8.83；系数 6.325、0.279、0.243、−0.207、0.500\\
y 方向异常值 & outlier\_toy6.csv & 第 3 点 $t_3=4.316$，$P=0.0497$；加入第 6 点后 $h_{66}=0.836$，$D_6=0.106$\\
偏相关 & score\_iq18.csv & 控制智商后数学与语文偏相关 0.4612；$t=2.013$，$P=0.062$\\
完全随机设计 & mice\_oneway30.csv & 设计矩阵秩为 3；\Verb|lm(x ~ g - 1)| 给出各组均数\\
随机区组 & rabbit\_block24.csv & 正规方程解与 \Verb|lm()| 相同；\Verb|anova()|\\
2×2 析因 & anemia\_factorial12.csv & SS = 1.6875、0.9075、0.3675\\
协方差分析 & chol\_ancova26.csv & 平行性 $P=0.935$；组别 $P=0.03755$（以 fit3 为分母 0.04202）；校正均数 5.491、6.386\\
单组重复测量 & weight\_rm10.csv & $W=0.1305$，一阶近似 $P=0.0882$，GG = 0.5397\\
多组重复测量 & anesthesia\_rm15.csv & $W=0.2931$，GG = 0.6787\\
\bottomrule
\end{tabular}
\end{table}

\section{常见报错与排查}
\label{app:errors}
\begingroup\small
\begin{longtable}{@{}p{5.2cm}p{4.2cm}p{4.8cm}@{}}
\caption{本课程代码里实际会遇到的报错}\\
\toprule
报错（中文环境的提示） & 原因 & 改法\\
\midrule
\endfirsthead
\toprule
报错 & 原因 & 改法\\
\midrule
\endhead
\bottomrule
\endfoot
\Verb|cannot open file 'd://data/a.csv'|（无法打开文件） & 课堂代码的绝对路径 & 用 .Rproj 打开，改成 \Verb|"data/datasetA.csv"|\\
\Verb|object 'lwt_group' not found|（找不到对象） & 变量在数据框里，没写 \Verb|data =|；或依赖 \Verb|attach()| & 函数里写 \Verb|data = da|\\
\Verb|could not find function "pivot_longer"| & 没有加载包 & \Verb|library(tidyr)|；包没装就先 \Verb|install.packages()|\\
\Verb|'newdata' had 1 row but variables found have 200 rows| & newdata 的列名与模型变量名不一致 & 列名照模型公式写\\
\Verb|contrasts can be applied only to factors with 2 or more levels| & 子集后某个因子只剩一个水平 & \Verb|droplevels()|，或检查筛选条件\\
\Verb|non-conformable arrays|（非整合陈列） & \Verb|tapply()| 结果是一维数组，和矩阵运算 & \Verb|as.vector()| 包一下\\
\Verb|argument "id" is missing|（缺少参数 “id”） & 对 regsubsets 对象调用 \Verb|vif()|、\Verb|coef()| & \Verb|vif()| 用在 lm 上；\Verb|coef(fit, id = 3)|\\
\Verb|Cannot add <ggproto> objects together| & ggplot 某行末尾漏了 \Verb|+|，或 \Verb|+| 写在了行首 & 每行末尾写 \Verb|+|\\
\Verb|The package "Hmisc" is required| & \Verb|mean_cl_normal| 依赖 Hmisc & 改用 \Verb|mean_se|，或安装 Hmisc\\
中文显示为乱码 & 文件编码 & \Verb|read.csv(..., fileEncoding = "UTF-8")|\\
因子水平变成 NA & \Verb|factor(x, levels = )| 中的名字与原值不一致（如 \Verb|"(110, Inf]"| 的空格） & 先 \Verb|levels(x)| 再复制\\
III 型结果与课件不同 & 没设和为零编码 & \Verb|options(contrasts = c("contr.sum", "contr.poly"))|\\
olsrr 阈值没生效 & 参数名写错被忽略 & 前进 / 后退用 \Verb|p_val|\\
\end{longtable}
\endgroup

\section{课堂原始代码与整理版对照}
\label{app:orig}
原件保存在 \Verb|课堂原始代码/|，没有改动；完整对照表在 \Verb|课堂原始代码/对照说明.md|。这里只列直接运行原件会出问题的地方：

\begin{itemize}
\item \textbf{所有原件}：路径 \Verb|d://data/…|（\Verb|a.csv| 即 datasetA，\Verb|hw.csv| 即 diabetes27\_hw）；大量 \Verb|attach()|。
\item \textbf{回归诊断.R}：\Verb|regsubsets(glu ~ ., data = d[,-1])| 把 age 也删掉了；\Verb|vif(fit)| 时 fit 已是 regsubsets 对象，报错；\Verb|kappa(scale(d))| 包含了因变量且是近似值。
\item \textbf{筛选自变量.R}：\Verb|predict(lm(y ~ x), …)| 中 x 未定义。
\item \textbf{stepwise.R}：第 17–18 行是粘贴的输出，source 时语法错误。
\item \textbf{第三讲.R}：\Verb|var.test(bwt ~ lwt_group)| 等没写 \Verb|data = da|；协方差分析 \Verb|anova()| 中分组在协变量前面。
\item \textbf{第四讲.R}：读入对象叫 \Verb|da|、后面用 \Verb|bp|；\Verb|pivot_longer| 前未加载 tidyr；\Verb|anova(lm(post ~ group + base))| 的分组一行未校正（第 \ref{sec:bp} 节）。
\end{itemize}

课件中与数据不符、本笔记已注明的地方：第二讲第 7 页残差 SD 的四等分数值；第四讲第 22 页 Mauchly 卡方近似的校正项；第四讲第 31 页折线图上的均数；第四讲第 10–11 页方差分析表中的“分组”一行。

\end{document}
